% main.tex -- draft paper for item 23 (Bolza surface, local extremality of lambda_1).
% Authoritative source: ../work/PROOF.md.  Every number below is copied from PROOF.md or from
% the certificate output files in ../work/ (see Appendix B).  Nothing was recomputed for this draft.
\documentclass[11pt]{amsart}

\usepackage[utf8]{inputenc}
\usepackage[T1]{fontenc}
\usepackage{lmodern}
\usepackage{amsmath,amssymb,amsthm}
\usepackage{mathtools}
\usepackage{booktabs}
\usepackage{array}
\usepackage{enumitem}
\usepackage[margin=1.1in]{geometry}
\usepackage{graphicx}
\usepackage[hidelinks]{hyperref}
\emergencystretch=3em

\newtheorem{theorem}{Theorem}[section]
\newtheorem{proposition}[theorem]{Proposition}
\newtheorem{lemma}[theorem]{Lemma}
\newtheorem{corollary}[theorem]{Corollary}
\theoremstyle{definition}
\newtheorem{definition}[theorem]{Definition}
\newtheorem{hypothesis}[theorem]{Hypothesis}
\theoremstyle{remark}
\newtheorem{remark}[theorem]{Remark}

\newtheorem*{maintheorem}{Main Theorem}

\newcommand{\R}{\mathbb{R}}
\newcommand{\Z}{\mathbb{Z}}
\newcommand{\C}{\mathbb{C}}
\newcommand{\Q}{\mathbb{Q}}
\newcommand{\Hyp}{\mathbb{H}^2}
\newcommand{\Teich}{\mathcal{T}_2}
\newcommand{\Oct}{\mathcal{C}}
\newcommand{\Sl}{\mathcal{S}}
\newcommand{\Sym}{\operatorname{Sym}}
\newcommand{\tr}{\operatorname{tr}}
\newcommand{\rank}{\operatorname{rank}}
\newcommand{\Isom}{\operatorname{Isom}}
\newcommand{\Spec}{\operatorname{Spec}}
\newcommand{\Ran}{\operatorname{Ran}}
\newcommand{\arccosh}{\operatorname{arccosh}}
\newcommand{\TT}{\mathrm{TT}}
\newcommand{\ls}{\lambda_*}
\newcommand{\gin}{g_{\mathrm{inv}}}
\newcommand{\smid}{s_{\mathrm{mid}}}
\newcommand{\svert}{s_{\mathrm{vert}}}
\newcommand{\xB}{x_B}
\newcommand{\wB}{w_B}
\newcommand{\vB}{v_B}
\newcommand{\FB}{F_B}
\newcommand{\FA}{F_A}

\begin{document}

\title[The Bolza surface is a strict local maximum of $\lambda_1$]%
{The Bolza surface is a strict local maximum of the first Laplace eigenvalue on genus-two Teichm\"uller space}

\author{Gordon Speagle}
\address{[Affiliation to be added]}
\email{gspeagle@gmail.com}

\date{Draft, \today}

\subjclass[2020]{58J50 (primary); 32G15, 35P15, 65G30, 65N30 (secondary)}
\keywords{Bolza surface, Laplace eigenvalues, hyperbolic surfaces, Teichm\"uller space,
eigenvalue optimisation, computer-assisted proof, interval arithmetic}

\begin{abstract}
Let $B$ be the Bolza surface, the most symmetric closed hyperbolic surface of genus two.
Strohmaier and Uski conjectured in 2013 that $B$ maximises the first positive Laplace
eigenvalue $\lambda_1$ among closed hyperbolic surfaces of genus two, and local maximality has
been asserted on numerical grounds. We give a computer-assisted proof that $B$ is a
\emph{strict local maximum} of $\lambda_1$ on the Teichm\"uller space $\Teich$. The proof does
not depend on the eigenvalue enclosure of Strohmaier and Uski: we certify
$\lambda_1(B)\ge3.8388757010$ ourselves by the Lehmann--Goerisch method with an
$H(\mathrm{div})$-conforming flux on the glued octagon. The proof is conditional on quoted results:
two lemmas and geodesic data of Fortier Bourque and Petri, the Selberg trace formula, and the
Lehmann--Goerisch theorem. Because $\lambda_1(B)$ has multiplicity three, $\lambda_1$ is not differentiable at $B$.
We analyse instead the analytic $3\times 3$ matrix $M$ that carries the eigenvalue cluster. Its
first-order part $A$ is traceless by symmetry and has rank exactly three. Its kernel $K$ is a
three-dimensional absolutely irreducible module for $\Isom(B)$, identified exactly as the
$(-1)$-eigenspace of the antipodal isometry. On $K$ the Hessian of the cluster sum is a scalar
multiple of an invariant inner product, and a certified Ky Fan bound along one explicit
deformation shows that this scalar is negative. The certified ingredients are a
Selberg-trace-formula count, conforming finite elements on the glued octagon in ball
arithmetic (Rayleigh--Ritz upper bounds and a Lehmann--Goerisch lower bound), exact character computations in $\Q(\zeta_{24})$, and certified rank computations.
The finite-dimensional logic of the proof has been checked in Lean~4, with the analytic and
computational inputs stated as explicit hypotheses.
\end{abstract}

\maketitle
\tableofcontents

%======================================================================
\section{Introduction}\label{sec:intro}
%======================================================================

\subsection{The problem}
For a closed hyperbolic surface $X$ let $0=\lambda_0<\lambda_1(X)\le\lambda_2(X)\le\cdots$ be the
spectrum of the Laplace--Beltrami operator. In genus two the \emph{Bolza surface} $B$, obtained
by gluing opposite sides of the regular hyperbolic octagon with angles $\pi/4$, is
distinguished in several ways. It has the largest isometry group (of order $96$) and the
longest systole~\cite{Str16,SU13}, and, conjecturally, the largest first eigenvalue.

\subsection{History}
Jenni~\cite{Jen84} estimated $\lambda_1(B)\in[3.83,3.85]$ and claimed that its multiplicity
is $3$; Cook~\cite{Coo18} found a gap in the multiplicity argument (both as reported
in~\cite{SU13,FBP-K}). Strohmaier and Uski~\cite[\S5]{SU13} computed
\[
\lambda_1(B)\approx 3.8388872588421995185866224504354645970819150157,
\]
stating that ``the rigorous implementation that uses very simple bounds in its current form
gives a rigorous error bound of $10^{-6}$''. They wrote: ``we conjecture that this is the
maximal value of the first eigenvalue of the Laplace operator for constant negative
curvature $-1$ in genus~2. This conjecture is supported by our calculations in Teichm\"uller
space (the corresponding analysis will be discussed elsewhere)''. We have not found the
promised analysis in print.

In his lecture notes, Strohmaier~\cite[pp.~22--23]{Str16} made the local statement explicit:
``Numerical evidence suggests that this is the global maximum for constant negative curvature
genus 2 surfaces. The reason for it being locally maximal is however its degeneracy. [\dots]
Numerically one can see that no matter what perturbation one chooses, none of the eigenvalues
$\lambda_1$, $\lambda_2$ and $\lambda_3$ has an extremal value at the Bolza surface.''
Cook~\cite[Conj.~5.1]{Coo18} restated the global conjecture and proposed a route through local
maximality: ``It is thought that by analyzing the nodal lines of eigenfunctions relating to
$\chi_8$ under perturbations in Teichm\"uller space, one would be able to show that the Bolza
surface is a local maximum for $\lambda_1$.''

Upper bounds now come close to $\lambda_1(B)$. Kravchuk, Maz\'a\v{c} and Pal~\cite[Thm.~1.1]{KMP}
proved, with exact rational verification, that every hyperbolic orbifold of genus $2$ satisfies
$\lambda_1<3.8388976481$. This is $1.04\times10^{-5}$ above $\lambda_1(B)$, and the authors
call it ``very strong evidence for the conjecture that the Bolza surface maximizes $\lambda_1$
among all $g=2$ orbifolds''. Bonifacio~\cite{Bon22a,Bon22b} obtained the numerical bootstrap
bound $\lambda_1\le 3.8388977$. Fortier Bourque and Petri proved that $\lambda_1(B)$ has
multiplicity $3$~\cite[Thm.~7.1]{FBP-K}, completing Cook's work, and developed linear
programming bounds for hyperbolic surfaces~\cite{FBP-LP}; they describe $B$ as the
``conjectured maximizer'' in genus~2. Radcliffe~\cite{Rad} gives numerical evidence that the
single-correlator bootstrap bound does not reach $\lambda_1(B)$ even in the limit
(``$\tilde\lambda_1^\infty-\lambda_1\ge\Delta_\infty\approx5.59683\times10^{-6}$''), so the
bootstrap route alone does not close the global problem as it stands. As of 2026 the global
conjecture is still described as open~\cite{HP,FBP-LP,Bon22b}; see also~\cite{PetriHDR,GPSX}.

To our knowledge, local maximality has not been proved anywhere. Since global maximality is
open, it cannot be deduced from a global result. The rigorous bounds~\cite{KMP,SU13} show only
that no genus-two surface exceeds $\lambda_1(B)$ by more than about $1.1\times10^{-5}$; they
give no local information.

\subsection{Main result}
We write $\ls:=\lambda_1(B)$, $E$ for the $\ls$-eigenspace, $G:=\Isom(B)$ and
$\Lambda:=\lambda_1+\lambda_2+\lambda_3$. The tangent space $T_B\Teich$ is identified with the
space of transverse-traceless (TT) tensors $\{\operatorname{Re}\varphi:\varphi\in Q(B)\}$, where
$Q(B)$ is the space of holomorphic quadratic differentials (Section~\ref{sec:setup}).

\begin{hypothesis}[SU; optional, not used in the proof]\label{hyp:SU}
The Bolza surface has a Laplace eigenvalue in the interval
$[3.83888625884,\,3.83888825884]$ \cite[\S5]{SU13}.
\end{hypothesis}
An earlier version of this proof assumed Hypothesis~\ref{hyp:SU} for its lower bound on
$\lambda_1(B)$. It is now replaced by the certified bound of Section~\ref{subsec:LG}, and
Hypothesis~\ref{hyp:SU} is used only as a cross-check and for alternative (sharper) constants,
which we record alongside.

\begin{maintheorem}
Assume the quoted results listed in Table~\ref{tab:deps}, namely~\cite[Lemmas~7.2--7.3]{FBP-K}, the
geodesic data of~\cite[\S7.4]{FBP-K}, the Selberg trace formula in the normalisation
of~\cite{FBP-K} with~\cite[Lemma~3.2]{FBP-K}, the Lehmann--Goerisch
theorem~\cite[Thms.~4.1, 4.3]{LiuPlum} (see also~\cite[Thm.~2.1]{BMPW} as reproduced in~\cite{Vej18}),
and the standard facts listed there. Then $B$ is a strict local maximum of $\lambda_1$ on $\Teich$:
there is a neighbourhood $N$ of $B$ in $\Teich$ such that $\lambda_1(X)<\lambda_1(B)$ for every
$X\in N\setminus\{B\}$. More precisely:
\begin{enumerate}[label=\textup{(\roman*)}]
\item $\lambda_1(B)$ has multiplicity $3$, and
$\lambda_1(B)\in[3.8388757010,\ 3.8388872830]$.
\item The first-order cluster map $A\colon T_B\Teich\to\Sym_0(E)$ has rank $3$. Its kernel $K$
is $3$-dimensional and is an absolutely irreducible $G$-module; it equals the
$(-1)$-eigenspace of the antipodal isometry.
\item Off $K$, $\lambda_1$ decreases to first order: $\lambda_1-\lambda_1(B)\le -c_A|w|/2$ in the
regime of case~(i) of Lemma~\ref{lem:criterion}. On $K$, the Hessian of $\Lambda$ satisfies
$D^2\Lambda|_K=c\cdot\gin$ for a $G$-invariant inner product $\gin$ and a constant $c<0$.
\item Along the explicit equiangular family $\FB$ (Section~\ref{subsec:FB}),
$(\Lambda\circ x)''(0)\le -9.3411$; the central value of the certified enclosure is $-14.274$.
\end{enumerate}
\end{maintheorem}

Both ends of the enclosure in~(i) are our own certificates: the upper end by conforming
Rayleigh--Ritz, the lower end by Lehmann--Goerisch (Section~\ref{subsec:LG}). The lower bound uses
only the trace-formula count $\lambda_4>4.0$ and the Lehmann--Goerisch theorem; it does not use
the multiplicity lemmas of~\cite{FBP-K}, which are still needed for multiplicity~$3$. The second-order
certificate needs $\lambda_1(B)\ge 3.8388029653$ and the first-order certificate
$\lambda_1(B)\ge3.83788$; the certified $3.8388757010$ exceeds both. If
Hypothesis~\ref{hyp:SU} is assumed in addition, the constant in~(iv) improves to $-12.8504$
(Remark~\ref{rem:SU}). The proof also uses published results that we quote rather than reprove
(the Selberg trace formula in the normalisation of~\cite{FBP-K}, the geodesic data and two lemmas
of~\cite{FBP-K} excluding low-dimensional representations in the $\lambda_1$-eigenspace, the
Lehmann--Goerisch theorem, the identification of $T_B\Teich$ with TT tensors, and the
Fischer--Marsden splitting); they are listed in Table~\ref{tab:deps}.

\subsection{Method}
The function $\lambda_1$ is not differentiable at $B$, because $\lambda_1(B)$ is triple. The
proof has four steps.

\emph{Cluster matrix.} Near $B$ the three lowest positive eigenvalues are the eigenvalues of a
real-analytic symmetric $3\times3$ matrix $M$ with $M(0)=\ls I$ (Theorem~\ref{thm:F1}). Write
$M(v)=\ls I+A(v)+Q_2(v)+O(|v|^3)$. The first-order term
$A(v)=[b_v(e_i,e_j)]$ is intrinsic on $T_B\Teich$ and $G$-equivariant (Lemma~\ref{lem:N}).

\emph{First order.} By symmetry $\tr A\equiv0$ (Proposition~\ref{prop:trA}). A certified
computation shows $A\ne0$ (Proposition~\ref{prop:Anonzero}). Exact representation theory then
forces $\ker A=K_0$, the $(-1)$-eigenspace of the antipodal isometry, with $\rank A=3$
(Proposition~\ref{prop:kerA}).

\emph{Second order.} A local-maximum criterion for degenerate clusters
(Lemma~\ref{lem:criterion}) reduces the theorem to negative definiteness of $D^2\Lambda$ on
$K$. Because $K$ is absolutely irreducible, $D^2\Lambda|_K=c\,\gin$ once $D^2\Lambda$ is known to
be $G$-invariant in a chart where $\Lambda$ is $C^2$ (Lemma~\ref{lem:Ginv}). So a single
direction in $K$ decides the sign of $c$. We take the tangent $\wB$ of the equiangular family
$\FB$, which lies in $K$ exactly (Lemma~\ref{lem:FBline}). A Ky Fan trace bound with a certified
trial space gives $(\Lambda\circ x)''(0)\le -9.3411$ (Proposition~\ref{prop:sign}).

\emph{Chart.} Six explicit trace functions are certified to be local coordinates of $\Teich$
at $B$ (Proposition~\ref{prop:G3}). This gives an analytic $6$-dimensional slice $\Sl$ of the
octagon space on which all of the above holds and which maps homeomorphically onto a
neighbourhood of $B$.

\subsection{What is computer-assisted}
The following steps are certified by computer. All are done in ball arithmetic (Arb, through
python-flint~0.9.0) or in exact rational arithmetic (sympy~1.14.0, exact arithmetic in
$\Q(\zeta_{24})$). Floating-point linear algebra is used only to \emph{propose} candidates
(test functions, eigenvectors, correctors), which are then verified rigorously.
\begin{enumerate}[label=(\alph*)]
\item Isolation of the cluster: a Selberg trace formula count (Section~\ref{subsec:LP}).
\item Upper bounds for $\lambda_3(B)$ by conforming finite elements on the glued octagon
(Section~\ref{subsec:FE}), and a lower bound for $\lambda_1(B)$ by the Lehmann--Goerisch method
with $H(\mathrm{div})$-conforming fluxes (Section~\ref{subsec:LG}).
\item $A\ne0$ (Section~\ref{subsec:Anonzero-cert}).
\item The second-order sign bound along $\FB$ (Section~\ref{subsec:sign-cert}).
\item The rank of the constraint Jacobian of the octagon space, and the rank of the
differential of the map to $\Teich$ (Section~\ref{subsec:rank-cert}).
\item The characters of $G$ on $T_B\Teich$ and the decomposition $T_B\Teich=K_0\oplus W_0$
(Section~\ref{sec:symmetry}; exact).
\end{enumerate}
Everything else is a written argument.

\subsection{Related work}
The pattern ``symmetry gives criticality, a computation gives non-degeneracy'' is familiar
from systolic geometry. Fortier Bourque, Mart\'inez-Granado and Vargas Pallete~\cite{FBMGVP}
proved that the extremal-length systole has a strict local maximum at $B$, combining eutaxy
(from symmetry) with perfection (from a rank computation); see also~\cite{FBR}. In our setting
the analogue of perfection fails: $\rank A=3<6=\dim\Teich$, so a genuine second-order
certificate on the kernel is needed. For flat tori, local maximality of eigenvalue functionals
is classical (see~\cite[\S3.3]{KLO}). For the spectral determinant, $B$ is a critical point by
symmetry~\cite{Kal26}, and numerically a local maximum~\cite[\S7.1]{SU13}. The area-normalised
problem over \emph{all} metrics is different: there the sharp genus-two constant $16\pi$ is
attained by a singular metric conformal to $B$~\cite{NS,JLNNP} (see also~\cite{KPS}), and the
second variation for that problem was only recently written down~\cite{Kar25}. For hyperbolic
surfaces that bound gives only $\lambda_1\le 4$~\cite[\S6]{FBP-K}.

Our result is consistent with Strohmaier's numerical assertion of local maximality, but it
qualifies the stated mechanism. Along directions in the $3$-dimensional kernel $K$ (for
example along $\FB$) all three eigenvalue branches are stationary to first order, so the
statement that ``none of the eigenvalues $\lambda_1,\lambda_2,\lambda_3$ has an extremal
value'' fails there; it holds on the open dense set of directions off $K$
(Remark~\ref{rem:Str16}).

\subsection{Organisation}
Section~\ref{sec:setup} sets up the octagon space and the pulled-back forms.
Section~\ref{sec:perturbation} proves analyticity of the cluster and the intrinsic nature of $A$.
Section~\ref{sec:symmetry} contains the exact representation theory.
Sections~\ref{sec:first} and~\ref{sec:second} carry out the first- and second-order analysis.
Section~\ref{sec:cert} describes the certified computations, including the lower bound for
$\lambda_1(B)$, and Section~\ref{sec:proof}
assembles the proof. Section~\ref{sec:lean} describes a Lean formalisation of the
finite-dimensional logic, and Section~\ref{sec:remarks} collects remarks and open problems.
Appendix~\ref{app:data} tabulates the certificate data and Appendix~\ref{app:repro} explains how
to reproduce them.

%======================================================================
\section{Setup}\label{sec:setup}
%======================================================================

\subsection{Hyperbolic plane}
We use the hyperboloid model $\Hyp=\{X\in\R^3:\mathsf B(X,X)=1,\ X_0>0\}$ with
$\mathsf B=\operatorname{diag}(1,-1,-1)$, and the associated Klein (projective) coordinates.
Orientation-preserving isometries form $SO^+(2,1)\cong PSL(2,\R)$.

\subsection{The octagon space}\label{subsec:octspace}
Let $x=(V_0,\dots,V_7)\in\R^{16}$ be the Klein coordinates of the vertices of a convex
hyperbolic octagon, with side $k$ joining $V_k$ and $V_{k+1}$ (indices mod $8$). The Bolza
gluing pairs side $k+4$ with side $k$ ($k=0,\dots,3$) by the orientation-preserving isometry
$g_k(x)\in SO^+(2,1)$ with $V_{k+4}\mapsto V_{k+1}$ and $V_{k+5}\mapsto V_k$. Explicitly,
$g_k=F(V_{k+1},V_k)\,F(V_{k+4},V_{k+5})^{-1}$, where $F(P,Q)=[P,e,n]$ is the
Lorentz-orthonormal frame with $e\propto Q-\mathsf B(P,Q)P$ and $n\propto \mathsf B(P\times e)$.

Let $\Oct\subset\R^{16}$ be the set of $x$ satisfying the five constraints
\[
L_{k+4}=L_k\ (k=0,\dots,3),\qquad \textstyle\sum_{k}\theta_k=2\pi,
\]
where $L_k$ are the side lengths and $\theta_k$ the interior angles. By the Poincar\'e polygon
theorem, each $x\in\Oct$ yields a closed hyperbolic surface $B_x$ of genus $2$, with the single
vertex cycle $0\to5\to2\to7\to4\to1\to6\to3\to0$. The point $\xB\in\Oct$ is the regular
octagon, with sides $2\arccosh(1+\sqrt2)=3.0571418389620$ and angles $\pi/4$; $B=B_{\xB}$.
By the certified rank of the constraint Jacobian (Proposition~\ref{prop:rankJ}) and the analytic
implicit function theorem, $\Oct$ is a real-analytic $11$-manifold near $\xB$, with
$T_{\xB}\Oct=\ker J$.

\subsection{The map to Teichm\"uller space}
The surface $B_x$ is marked through the isomorphism $\Gamma_{\xB}\to\Gamma_x$,
$g_k(\xB)\mapsto g_k(x)$; this marking is induced by the mesh homeomorphism $\Phi_x$ below and is
constant in $x$ up to isotopy. This defines $\pi\colon\Oct\to\Teich$. The map $\pi$ is
continuous, because it factors through the analytic holonomy map $x\mapsto(g_k(x))_k$ and the
quotient by conjugation.

\subsection{The mesh map and the pulled-back forms}\label{subsec:mesh}
The octagon is triangulated by a fixed fan of geodesic triangles. Nodes on the sides are placed
at exact hyperbolic-arclength fractions $j/n$ and matched $j\leftrightarrow n-j$ by the side
pairings. Interior nodes are given by explicit isometry-equivariant formulas: the centre is the
normalised vertex sum, and interpolation uses the symmetric hyperbolic-arclength formula
$\sinh(sL)/\sinh L$. The parameter $n$ is the mesh resolution.

On a triangle with hyperboloid vertices $P_0,P_1,P_2$ set
\[
W(a,b)=aP_0+bP_1+(1-a-b)P_2,\qquad X(a,b)=W/\sqrt{\mathsf B(W,W)}.
\]
Cellwise, these maps define a piecewise-analytic, bi-Lipschitz homeomorphism
$\Phi_x\colon B\to B_x$ that is compatible with the gluings. Pulling back the Dirichlet and
$L^2$ forms of $B_x$ gives, on $H^1(B)$,
\begin{equation}\label{eq:forms}
q_x(f)=\sum_T\int(\nabla f)^{\mathsf T}a_T\nabla f,\qquad
m_x(f)=\sum_T\int\rho_T f^2,
\end{equation}
with
\[
a_T=\frac{\operatorname{adj}(LL^{\mathsf T}-QG)}{D\sqrt Q},\qquad \rho_T=D\,Q^{-3/2},
\]
where $Q=\mathsf B(W,W)\ge1$, $W_1=P_0-P_2$, $W_2=P_1-P_2$, $L_i=\mathsf B(W_i,W)$,
$G_{ij}=\mathsf B(W_i,W_j)$ and $D=|\det[P_0P_1P_2]|$. (These formulas were derived by hand and
checked against a finite-difference pullback; the certified area enclosure $\to4\pi$ is a
further check.) The pencil $(q_x,m_x)$ has the Laplace spectrum of $B_x$.

\subsection{Notation}\label{subsec:notation}
\begin{itemize}
\item $E$ is the $\ls$-eigenspace of $B$; $e_1,e_2,e_3$ is an $m_0$-orthonormal basis.
\item $G=\Isom(B)$ has order $96$; $G^+=\Isom^+(B)\cong GL(2,3)$; $J=r^4$ is the hyperelliptic
involution.
\item $T_B:=T_B\Teich\cong\{\text{TT tensors on }(B,h_0)\}=\{\operatorname{Re}\varphi:\varphi\in
Q(B)\}$, where $h_0$ is the hyperbolic metric of $B$
(\cite[\S1]{Maz}, citing~\cite{Tro92}; see also~\cite[\S2]{DW}). As a real $G$-module,
$G^+$ acts $\C$-linearly by pull-back and orientation-reversing elements act $\C$-antilinearly.
\item Octagon symmetries: $r$ is the rotation by $\pi/4$ (relabelling $k\mapsto k+1$), acting
on $\R^{16}$ by $(rx)_k=\mathrm{Rot}_{\pi/4}(x_{k-1})$; $\smid$ is the reflection in an axis
through two opposite side midpoints; $\svert$ is the reflection in a diagonal through two
opposite vertices.
\item In the model $y^2=x^5-x$ (Section~\ref{sec:symmetry}): $t$ is an element of order $3$,
$u\colon x\mapsto i/x$, $\kappa\colon(x,y)\mapsto(\bar x,\bar y)$, and $a$ is the antipodal
isometry, the order-$4$ orientation-reversing lift of $x\mapsto-1/\bar x$.
\end{itemize}

%======================================================================
\section{Perturbation of the eigenvalue cluster}\label{sec:perturbation}
%======================================================================

\subsection{Analyticity of the cluster}
A \emph{pulled-back family} is a map $t\mapsto(a_t,\rho_t)\in L^\infty$, $t\in U\subset\R^d$, of
measurable coefficient fields on $B$ (written in the reference simplices of the mesh), defining
forms $q_t(f,g)=\int a_t\nabla f\cdot\nabla g$ and $m_t(f,g)=\int\rho_tfg$ on $H:=H^1(B)$. We assume:
\begin{enumerate}[label=(H\arabic*)]
\item\label{H1} $t\mapsto(a_t,\rho_t)\in L^\infty$ is $C^k$, $k\in\{2,3,\dots,\infty,\omega\}$; for
$k=\omega$ it extends holomorphically to a complex neighbourhood of $0$;
\item\label{H2} $a_0\ge\alpha I$ and $\rho_0\ge\alpha>0$ almost everywhere;
\item\label{H3} for real $t$, $(q_t,m_t)$ is the Dirichlet/$L^2$ pair of a closed hyperbolic
surface $B_t$ pulled back by a bi-Lipschitz homeomorphism;
\item\label{H4} $\Spec(B)\cap[\ls-r_0,\ls+r_0]=\{\ls\}$ with multiplicity $3$.
\end{enumerate}
The octagon family on $\Oct$ satisfies \ref{H1}--\ref{H3} with $k=\omega$: $a_T$ and $\rho_T$ are
analytic in the vertex coordinates, at $\xB$ we have $Q\ge1$ and $D>0$ uniformly over the
finitely many elements, and the interpolation $\sinh(sL)/\sinh L$ is analytic in $\cosh L$. The
same holds on every analytic sub-family. Hypothesis~\ref{H4} holds with $r_0=0.16$ by
Lemma~\ref{lem:isolation}.

\begin{theorem}[Cluster reduction]\label{thm:F1}
Assume \ref{H1}--\ref{H4} and let $\Gamma=\{|z-\ls|=r_0/2\}$. For $t$ in a neighbourhood of $0$:
\begin{enumerate}
\item exactly $\lambda_1(t),\lambda_2(t),\lambda_3(t)$ lie inside $\Gamma$, and none lies on $\Gamma$;
\item the Riesz projection $P(t)=(2\pi i)^{-1}\oint_\Gamma(zM_t-Q_t)^{-1}M_t\,dz$ is $C^k$
(holomorphic if $k=\omega$), and for real $t$ it is the $m_t$-orthogonal projection onto the
cluster;
\item there is a $C^k$ symmetric matrix $M(t)\in\Sym(3)$ whose eigenvalues are
$\lambda_1(t),\lambda_2(t),\lambda_3(t)$; in particular $\Lambda=\tr M$ is $C^k$
(real-analytic if $k=\omega$);
\item $M(0)=\ls I$ and $dM(0)[v]=A(v):=[b_v(e_i,e_j)]_{ij}$, where
$b_v=\partial_vq-\ls\partial_vm$;
\item $M(v)=\ls I+A(v)+Q_2(v)+R(v)$ with $Q_2$ quadratic and $|R(v)|\le C|v|^3$ ($k\ge3$) or
$R(v)=o(|v|^2)$ ($k=2$), uniformly on a ball.
\end{enumerate}
\end{theorem}

\begin{proof}
\emph{Step 1.} Let $Q_t,M_t\colon H_\C\to H_\C'$ be the operators of $q_t,m_t$ (bilinear pairing).
The maps $a\mapsto Q[a]$, $\rho\mapsto M[\rho]$ are bounded linear from $L^\infty$, so
$t\mapsto Q_t,M_t$ inherit the regularity of~\ref{H1}. $M_t$ is compact, since it factors through
$L^2$ (Rellich).

\emph{Step 2.} For real $t$ near $0$, $q_t+m_t$ is an inner product equivalent to the $H^1$
norm, and $S_t=(Q_t+M_t)^{-1}M_t$ is compact, self-adjoint and positive for it. The spectral
theorem gives an $m_t$-orthonormal eigenbasis $(\varphi_k)$ with eigenvalues $\lambda_k(t)\to\infty$,
which by~\ref{H3} are the eigenvalues of $B_t$, and
$(zM_t-Q_t)^{-1}M_tf=\sum_km_t(f,\varphi_k)(z-\lambda_k)^{-1}\varphi_k$ for $z\notin\{\lambda_k\}$.

\emph{Step 3.} By~\ref{H4}, $zM_0-Q_0$ is invertible for $z\in\Gamma$ with uniformly bounded
inverse. By the Neumann series, applied uniformly on the compact contour, $zM_t-Q_t$ is
invertible for $(t,z)\in U'\times\Gamma$, and the resolvent is $C^k$ in $t$ uniformly in $z$.

\emph{Step 4.} Differentiating under the contour integral shows that $P(t)$ is $C^k$. For real $t$
the expansion of Step~2 identifies $P(t)$ with the $m_t$-orthogonal projection onto the
eigenvectors inside $\Gamma$. Idempotents at distance $<1$ have equal rank, so
$\rank P(t)=3$. Applying the same argument to a contour $\Gamma'$ enclosing
$[-\tfrac12,\ls+r_0/2]$, which at $t=0$ contains exactly $\{0,\ls,\ls,\ls\}$ (here we use
$\Spec(B)\cap(0,3.3]=\emptyset$), shows that exactly four eigenvalues lie below $\ls+r_0/2$. One
is $0$, so the three inside $\Gamma$ are $\lambda_1,\lambda_2,\lambda_3$.

\emph{Step 5.} Put $f_i(t)=P(t)e_i$, $G(t)_{ij}=m_t(f_i,f_j)$ (so $G(0)=I$),
$\tilde f=f\,G^{-1/2}$ and $M(t)_{ij}=q_t(\tilde f_i,\tilde f_j)$. Then $M$ is $C^k$ and symmetric,
and for real $t$, $(\tilde f_j)$ is an $m_t$-orthonormal basis of $\Ran P(t)$, so the
eigenvalues of $M(t)$ are $\lambda_1,\lambda_2,\lambda_3$.

\emph{Step 6.} Since $\tilde f_i(0)=e_i$ and $q_0(g,e_j)=\ls m_0(g,e_j)$, differentiating
$M_{ij}$ and the identity $m_t(\tilde f_i,\tilde f_j)=\delta_{ij}$ gives
$M'_{ij}=q'(e_i,e_j)-\ls m'(e_i,e_j)=b_v(e_i,e_j)$. Item~5 is Taylor's theorem.
\end{proof}

\begin{remark}
Only the individual eigenvalues can fail to be differentiable in several parameters (compare the
eigenvalues $\pm|t|$ of $\left(\begin{smallmatrix}t_1&t_2\\t_2&-t_1\end{smallmatrix}\right)$); the
total projection and every symmetric function of the cluster remain analytic. The same
Riesz-projection argument for $t\in\R^N$ appears in~\cite[Prop.~2.5]{Cai26} for a family with a
fixed Dirichlet form and a multiplicative weight, which is an analogue rather than our
variable-coefficient pencil. For the one-parameter theory see Kato~\cite[Ch.~VII]{Kato}
(as quoted in~\cite[\S4.1]{Dur}).
\end{remark}

\subsection{The first-order map is intrinsic and equivariant}

For an $L^\infty$ symmetric $2$-tensor $\dot h$ on $B$ set
$b_{\dot h}(e,e'):=\frac{d}{ds}\big|_{0}(q-\ls m)_{h_0+s\dot h}(e,e')$ for $e,e'\in E$.

\begin{lemma}\label{lem:N}
\begin{enumerate}[label=\textup{(\alph*)}]
\item For every $v\in T_{\xB}\Oct$ there is a TT tensor $\sigma(v)$, depending linearly and
only on $v$, such that $A_{\mathrm{oct}}(v):=[b_v(e_i,e_j)]=A(\sigma(v)):=[b_{\sigma(v)}(e_i,e_j)]$.
Thus $A$ is well defined on $T_B$.
\item $A(g\cdot\sigma)=\rho_E(g)A(\sigma)\rho_E(g)^{\mathsf T}$ for all $g\in G$, of either orientation.
\item For an octagon symmetry $\gamma$ acting linearly on $\Oct$, $\sigma(\gamma v)=\gamma_B\cdot\sigma(v)$.
\end{enumerate}
\end{lemma}

\begin{proof}
\emph{Step 1 (first variation).} Using $d(h^{ij})=-h^{ik}\dot h_{kl}h^{lj}$ and
$d\sqrt h=\tfrac12(\tr\dot h)\sqrt h$,
\[
b_{\dot h}(e,e')=-\int_B\langle\dot h,T(e,e')\rangle\,dA,\qquad
T(e,e'):=\Sym(de\otimes de')-\tfrac12\big(\langle de,de'\rangle-\ls ee'\big)h_0 .
\]
The formula is tensorial and may be evaluated cellwise; it agrees with the $b_v$ of
Theorem~\ref{thm:F1}, since the coefficients there are $a=\sqrt h\,h^{-1}$ and $\rho=\sqrt h$ and
$t\mapsto h_t$ is $C^1$ into $L^\infty$. (Compare~\cite[Lemma~2.1]{ESI}.)

\emph{Step 2 ($T$ is divergence-free on $E$).} Eigenfunctions are smooth and
$\tr\nabla^2e=-\ls e$. A direct computation gives
\[
2\operatorname{div}T=(\tr\nabla^2e)\,de'+(\tr\nabla^2e')\,de+\nabla^2e(\nabla e',\cdot)+\nabla^2e'(\nabla e,\cdot)
-d\langle de,de'\rangle+\ls(e\,de'+e'\,de)=0 .
\]

\emph{Step 3 (Lie derivatives drop out).} For a $W^{1,\infty}$ vector field $X$,
$L_Xh_0=2\Sym\nabla X^\flat\in L^\infty$, and integration by parts against the smooth $T$ gives
$b_{L_Xh_0}=2\int\langle X,\operatorname{div}T\rangle=0$.

\emph{Step 4 (octagon tangents are TT plus a Lie derivative).} Let $x(t)$ be a $C^1$ curve in
$\Oct$ through $\xB$, $h_t=\Phi_t^*h_{B_t}$, and let $D_t\colon\Hyp\to\Hyp$ be the developing map
of $h_t$, which is $\rho_t$-equivariant, cellwise analytic in $(t,p)$, and normalised by $D_0=\mathrm{id}$.
Fix a smooth radial kernel $\eta_\varepsilon$ supported in an $\varepsilon$-ball and define the
Riemannian centre of mass
\[
\tilde\Psi_t(p):=\operatorname*{argmin}_{q\in\Hyp}\int d(q,D_t(y))^2\,\eta_\varepsilon(d(p,y))\,dA(y).
\]
The functional is strictly convex and proper, so $\tilde\Psi_t$ is well defined; it is
$\rho_t$-equivariant because the kernel is deck-invariant; it is jointly $C^\infty$ in $(t,p)$ by
the implicit function theorem (the Hessian in $q$ is positive definite, $p$-derivatives fall on
the kernel and $t$-derivatives on $D_t$ inside the integral); $\tilde\Psi_0=\mathrm{id}$; and
$d\tilde\Psi_t=I+O(t)$, so $\tilde\Psi_t$ is a diffeomorphism for small $|t|$. It descends to
$\Psi_t\colon B\to B_t$. Then $k_t:=\Psi_t^*h_{B_t}$ is a smooth family of smooth hyperbolic metrics
with $k_0=h_0$, and $\iota_t:=\Psi_t^{-1}\circ\Phi_t$ is an isometry $(B,h_t)\to(B,k_t)$ isotopic to
the identity. The mixed derivative $\partial_t\nabla_p\iota_t$ exists and is continuous on each
closed cell, so $X:=\partial_t\iota_t|_0$ is Lipschitz and $\dot h=\dot k+L_Xh_0$ almost everywhere.

By the Fischer--Marsden splitting, as stated in~\cite[Rem.~3.3]{Maz}, every symmetric tensor
$S$ decomposes orthogonally as $S=S_{tt}+L_Xg+((-\Delta_gf+f)g+g\nabla^2f)$, where
$S_{tt}+L_Xg\in\ker dK_g$ and the last term lies in $(\ker dK_g)^\perp$. Applied to
$S=\dot k\in\ker dK_{h_0}$, the last term vanishes, so $\dot k=\sigma+L_Yh_0$ with $\sigma$ TT and $Y$
smooth by elliptic regularity. Set $\sigma(v):=\sigma$. Steps~1 and~3 give
$A_{\mathrm{oct}}(v)=A(\sigma(v))$. Two choices of $\Psi_t$ differ by $\mathrm{Diff}_0$, which changes $\dot
k$ by a Lie derivative, so $\sigma(v)$ is independent of the choices; it is linear in $v$.

\emph{Step 5 (equivariance).} Let $g$ be an isometry of $h_0$, acting on TT tensors by
$(g^{-1})^*$ and on $E$ by $\rho_E(g)e=e\circ g^{-1}$. Then
$T(e\circ g^{-1},e'\circ g^{-1})=(g^{-1})^*T(e,e')$ and $dA$ is $g$-invariant, which gives~(b).
On $\operatorname{Re}Q(B)$ this is $\operatorname{Re}((g^{-1})^*\varphi)$ for holomorphic $g$ and
$\operatorname{Re}(\overline{(g^{-1})^*\varphi})$ for antiholomorphic $g$.

\emph{Relabelling.} The fan mesh is equivariant under octagon symmetries, so
$\Phi_{\gamma x}=\Gamma\circ\Phi_x\circ\gamma_B^{-1}$ and $h_{\gamma x(t)}=(\gamma_B^{-1})^*h_{x(t)}$. Since
the TT projection is natural under isometries of $h_0$, this gives~(c).
\end{proof}

%======================================================================
\section{Symmetry: characters and the $K_0$ reduction}\label{sec:symmetry}
%======================================================================

\subsection{The algebraic model}
By~\cite[\S7.1]{FBP-K}, $\Isom^+(B)/\langle r^4\rangle$ is the octahedral rotation group acting on
$B/\langle J\rangle\cong\hat\C$ and permuting the six Weierstrass images; $\Isom^+(B)\cong GL(2,3)$
and $\Isom(B)\cong GL(2,3)\rtimes\Z/2$. Since a finite $S_4\subset PSL(2,\C)$ is conjugate to the
standard octahedral group, whose unique orbit of size $6$ is $\{0,\infty,\pm1,\pm i\}$, we may
take $B\cong\{y^2=x^5-x\}$.

The maps
\[
r\colon(x,y)\mapsto(ix,\zeta_8y),\qquad t\colon x\mapsto\frac{x+i}{x-i}\ (\text{with }y\text{-factor }2+2i),\qquad
\kappa\colon(x,y)\mapsto(\bar x,\bar y)
\]
generate a group of order $96$, of which $48$ elements are holomorphic and form a copy of
$GL(2,3)$ (element orders $1{:}1$, $2{:}13$, $3{:}8$, $4{:}6$, $6{:}8$, $8{:}12$). This was verified
in exact arithmetic over $\Q(\zeta_{24})$, including the curve identity for every element. Since
$|\Isom(B)|=96$~\cite{FBP-K}, this is all of $G$.

The space $Q(B)=\Sym^2H^0(K_B)$ has basis $(dx/y)^2$, $x(dx/y)^2$, $x^2(dx/y)^2$. An automorphism
with M\"obius matrix $\left(\begin{smallmatrix}a&b\\c&d\end{smallmatrix}\right)$ and $y$-factor $\mu$ pulls back
$(dx/y,\,x\,dx/y)$ by $P=(\det/\mu)\left(\begin{smallmatrix}d&b\\c&a\end{smallmatrix}\right)$; for
holomorphic $g$ the real character on $T_B\cong Q(B)\cong\R^6$ is $\chi_T(g)=2\operatorname{Re}\tr\Sym^2P$.
For orientation-reversing $g$ the action is $\C$-antilinear, and an antilinear map has real trace
$0$.

\begin{proposition}[Characters; exact]\label{prop:chars}
On $T_B$:
$\chi_T(r)=\chi_T(r^3)=\chi_T(r^2)=-2$, $\chi_T(J)=6$, $\chi_T(t)=0$, and $\chi_T(g)=0$ for every
orientation-reversing $g$. Hence $\langle\chi_T,\chi_T\rangle=2$ and $\langle\chi_T,1\rangle=0$; in
particular $T_B^G=0$.
\end{proposition}

These values were also reproduced with an independently written pull-back model.

\subsection{The antipodal splitting}
The antipodal isometry $a$ has conjugacy class $\{a,a^3=aJ\}$ of size $2$, so $a$ is central modulo
$J$, and $a^2=J$. Since $\chi_T(J)=6=\dim T_B$, $J$ acts trivially on $T_B$, so $a^2=\mathrm{id}$ on
$T_B$ and
\[
T_B=W_0\oplus K_0,\qquad W_0:=\{a=+1\},\quad K_0:=\{a=-1\},
\]
both $G$-invariant. Exactly, both have dimension $3$, $\langle\chi_\pm,\chi_\pm\rangle=1$ and
$\langle\chi_+,\chi_-\rangle=0$. Hence $K_0$ and $W_0$ are absolutely irreducible and
non-isomorphic. (In the labelling of Cook's character table~\cite[App.~A]{Coo18},
$K_0\cong\chi_7$ and $W_0\cong\chi_{10}$; this is not used.) On the octagon elements,
$\chi_{K_0}(r)=-1$, $\chi_{K_0}(\smid)=1$, $\chi_{K_0}(\svert)=-1$, and
$\chi_{W_0}(r)=-1$, $\chi_{W_0}(\smid)=-1$, $\chi_{W_0}(\svert)=1$; both vanish on $t$.

\begin{lemma}[Isolation of $E$; see Lemma~\ref{lem:isolation}]
$\dim E=3$ and $E$ is an irreducible real $G$-module.
\end{lemma}

\begin{proposition}\label{prop:K0kerA}
$K_0\subset\ker A$.
\end{proposition}

\begin{proof}
$E$ is irreducible of odd real dimension, so $\operatorname{End}_G(E)=\R$ and the central element $J$
acts on $E$ by $\pm1$. If $J=-1$ then $\rho_E(a)^2=-I$ on $\R^3$, which is impossible
($\det^2=-1$). So $J$ acts trivially, $\rho_E(a)$ commutes with $\rho_E(G)$ (as $a$ is central
modulo $J$), and $\rho_E(a)=\pm I$. By Lemma~\ref{lem:N}(b), $A(a\cdot v)=A(v)$. For $v\in K_0$,
$a\cdot v=-v$, so $A(v)=0$.
\end{proof}

\begin{proposition}\label{prop:kerA}
If $A\ne0$, then $\ker A=K_0$ and $\rank A=3$.
\end{proposition}

\begin{proof}
$A|_{W_0}\colon W_0\to\Sym(E)$ is $G$-equivariant and $W_0$ is irreducible, so $A|_{W_0}$ is $0$ or
injective. Since $A\ne0$ and $A|_{K_0}=0$, it is injective.
\end{proof}

\begin{proposition}\label{prop:scalar}
Every $G$-invariant symmetric bilinear form on $K_0$ is $c\cdot\gin$ for a fixed $G$-invariant
inner product $\gin$ and some $c\in\R$.
\end{proposition}

\begin{proof}
$K_0$ is absolutely irreducible (Schur).
\end{proof}

\begin{remark}[Consistency, not used]
The full Clifford--Gallagher computation shows $\Sym_0^2E\cong D\oplus W_0$ with $D$
two-dimensional irreducible, for each of the four candidate three-dimensional irreducibles. A
certified computation gives $\tr(g|E)\in[1\pm0.0352]$ for $g=r,\smid,\svert$ at $n=96$, so $E$ is
the octahedral representation, in agreement with~\cite[Thm.~3.9]{Coo18}. Neither fact is used.
\end{remark}

\subsection{Why the octagon group is not enough}\label{subsec:D8}
Restricted to the octagon group $D_8$ (16 elements), $K_0$ has norm
$(9+9+4\cdot1+2\cdot1+4\cdot1+4\cdot1)/16=2$, so $K_0|_{D_8}=\varepsilon_s\oplus(\text{2-dim})$.
$D_8$-invariant quadratic forms on $K_0$ thus form a two-parameter family, and one certified
direction controls only the $\varepsilon_s$-line. The non-octagon element $t$ is genuinely needed,
both for irreducibility of $W_0$ and for Proposition~\ref{prop:scalar}; the corresponding
invariance of the Hessian is proved in Lemma~\ref{lem:Ginv}.

\subsection{The equiangular family}\label{subsec:FB}
Let $\FB(p)$ be the octagon with all vertices at hyperbolic radius $R(p)$, $V_{2j}$ at angle
$j\pi/2$ and $V_{2j+1}$ at angle $\varphi+j\pi/2$, where $p=\varphi-\pi/4$. The $\pi/2$-rotation
and the reflection $\theta\mapsto\varphi-\theta$ act transitively on the vertices, so all angles are
equal, and $L_{k+4}=L_k$ by the $\pi$-rotation. Splitting the angle at $V_0$ along the spokes, the
surface condition (angle $=\pi/4$) is
\[
C^2xy-C(x+y)-1=0,\qquad C=\cosh R,\ x=\tan(\varphi/2),\ y=\tan(\pi/4-\varphi/2),
\]
so that $C=3+2\sqrt2$ at $p=0$. (The Taylor jets of the constraints along this curve were
certified to vanish to second order in Arb.) $\FB(p)$ is invariant under $r^2$ and $\smid$, and
$r$ maps $\FB(p)$ to $\FB(-p)$. Write $x(p)=\FB(p)$, $\vB=x'(0)$ and $\wB:=\sigma(\vB)$.

\begin{lemma}\label{lem:FBline}
$\wB\in K_0\setminus\{0\}$.
\end{lemma}

\begin{proof}
By exact linear algebra, $\{v\in T_B:r v=-v,\ \smid v=v\}$ is a line, and its spanning vector
satisfies $a v=-v$, so the line lies in $K_0$. By Lemma~\ref{lem:N}(c), $r\wB=-\wB$ and
$\smid\wB=\wB$, so $\wB$ lies on this line.

To see $\wB\ne0$: the pairing $g_0$ of $\FB(p)$ is the translation along the bisector of side $0$,
and hyperbolic trigonometry gives
$d\ell(g_0)/dp=-\cos(\pi/4)/(\sin^2(\pi/8)\sinh r_{\mathrm{in}})=-\sqrt{2+2\sqrt2}=-2.19736822693562$
and $d\ell(g_1)/dp=+2.19736822693562$, where $r_{\mathrm{in}}$ is the inradius of the regular octagon.
Thus $\ell(g_0)-\ell(g_1)$, a function on $\Teich$, has nonzero derivative along $\pi\circ\FB$.
By Lemma~\ref{lem:chart}(b) below, $d(\tau\circ x)(\vB)=d\tau(\wB)$ for length-type functions $\tau$,
so $\wB\neq0$. (The certificate of Proposition~\ref{prop:sign} also forces $\wB\ne0$.)
\end{proof}

\begin{remark}
The reflection type matters: for $\svert$-type reflections the line
$\{rv=-v,\ sv=v\}$ lies in $W_0$, not in $K_0$. So the identification of the symmetry of $\FB$ as
an $\smid$-type reflection (fixing four Weierstrass points) is load-bearing. Here $\smid$ fixes
the centre, the vertex point, the midpoint class on its axis and one further midpoint class via
the gluing; $\svert$ fixes only the centre and the vertex point.
\end{remark}

%======================================================================
\section{First-order analysis}\label{sec:first}
%======================================================================

\begin{proposition}\label{prop:trA}
$\tr A\equiv0$; equivalently $d\Lambda(\xB)=0$ on $T_{\xB}\Oct$.
\end{proposition}

\begin{proof}[First proof (octagon coordinates)]
The map $r\colon\R^{16}\to\R^{16}$ is linear, maps $\Oct$ to $\Oct$, and $B_{rx}$ is isometric to
$B_x$, so $\Lambda\circ r=\Lambda$. Hence $d\Lambda=d\Lambda\circ P$ on $T:=T_{\xB}\Oct$, with
$P=\frac18\sum_{j<8}r^j$ the projection onto $T^r$. The fixed space $(\R^{16})^r=\{\delta_k=\mathrm{Rot}^k\delta_0\}$
is two-dimensional (radial and tangential). The radial vector is not in $T$, because
$d(\text{angle sum})/dR=-5.5729711960\ne0$ (certified; Proposition~\ref{prop:rankJ}; it equals the
closed form from $\cot(\theta/2)=\cosh R\tan(\pi/8)$). The tangential vector is the infinitesimal
rotation, along which $\Lambda$ is constant. So $T^r$ is spanned by the rotation and
$d\Lambda=0$. By Theorem~\ref{thm:F1}(4), $d\Lambda(v)=\tr A(v)$.
\end{proof}

\begin{proof}[Second proof (Schur)]
$\tr A$ is a $G$-invariant linear functional on $T_B$ (Lemma~\ref{lem:N}(b)), and
$\langle\chi_T,1\rangle=0$ (Proposition~\ref{prop:chars}).
\end{proof}

In fact $G$-equivariance under the single element $r$ suffices, since $r$ has no fixed vector on
$T_B$; this is how the statement is formalised in Section~\ref{sec:lean}.

\begin{proposition}[$A\ne0$; certified]\label{prop:Anonzero}
Let $\delta\in T_{\xB}\Oct$ be the alternating-radius tangent: vertex radii
$R+t(-1)^k+\varsigma(t)$, angles $k\pi/4$, with $\varsigma(t)$ restoring the angle sum and
$\varsigma'(0)=0$. Then there is an exact eigenfunction $e\in E$ with
\[
b_\delta(e,e)\le-3.0383476242423058\cdot m(u)<0
\]
($n=192$; at $n=96$ the bound is $-0.9989201873540878\cdot m(u)$; under Hypothesis~\ref{hyp:SU} the
bounds are $-3.0875694043934345\cdot m(u)$ and $-1.0282409843057394\cdot m(u)$), where $u$ is the finite-element
function of Section~\ref{subsec:Anonzero-cert} and $e=P_Eu$. Consequently $A(\sigma(\delta))\ne0$.
\end{proposition}

The deformation $\delta$ preserves $L_{k+4}=L_k$ by the $\pi$-rotation symmetry; $\varsigma$ is
analytic by the implicit function theorem since $\partial(\sum\theta)/\partial\varsigma=-5.573\ne0$, and
$\varsigma'(0)=0$ because all $\partial\sum\theta/\partial R_k$ are equal at $\xB$ while
$\sum(-1)^k=0$. Up to scale $\delta$ is the tangent of the equilateral family $\FA$. The proof is
in Section~\ref{subsec:Anonzero-cert}. Identifying $[b_\delta(e_i,e_j)]$ with the first-order
variation along $\delta$ only needs the one-parameter case of Theorem~\ref{thm:F1}.

\begin{corollary}\label{cor:rank}
$K:=\ker A=K_0$, and $\rank A=3$.
\end{corollary}

\begin{proof}
Propositions~\ref{prop:K0kerA}, \ref{prop:kerA} and~\ref{prop:Anonzero}.
\end{proof}

%======================================================================
\section{Second-order analysis}\label{sec:second}
%======================================================================

\subsection{A local-maximum criterion for a degenerate cluster}

\begin{lemma}\label{lem:criterion}
Let $M(v)$ be a symmetric $3\times3$ matrix function on a neighbourhood of $0\in\R^6$ with
$M(v)=\ls I+A(v)+Q_2(v)+R(v)$, where $A$ is linear, $Q_2$ is quadratic and $|R(v)|=o(|v|^2)$.
Let $K=\ker A$. Assume
\begin{enumerate}[label=\textup{(\alph*)}]
\item $\tr A\equiv0$;
\item $\tr Q_2(k)\le-\kappa|k|^2$ for all $k\in K$, for some $\kappa>0$.
\end{enumerate}
Then $\lambda_{\min}M(v)<\ls$ for all sufficiently small $v\ne0$.
\end{lemma}

\begin{proof}
Let $W$ be a complement of $K$; $A|_W$ is injective. Each $A(w)$, $w\ne0$, is traceless and nonzero
and so has a negative eigenvalue; by compactness
$c_A:=\min_{w\in W,|w|=1}(-\lambda_{\min}A(w))>0$. This is the first use of~(a). Write $v=w+k$
with $w\in W$, $k\in K$, and choose $C''$ with $\|Q_2(v)\|+|R(v)|\le C''(|w|^2+|k|^2)$ on a small ball.

\emph{Case (i): $|w|\ge C|k|^2$, with $C>4C''/c_A$.} Then $w\ne0$, and by Weyl's inequality
\[
\lambda_{\min}M(v)-\ls\le-c_A|w|+C''|w|^2+C''|k|^2\le-c_A|w|+C''|w|^2+\tfrac{c_A}{4}|w|<-\tfrac{c_A}{2}|w|+C''|w|^2,
\]
which is negative for small $v$.

\emph{Case (ii): $|w|<C|k|^2$.} Then $k\ne0$ and $|v|\asymp|k|$. Using $\tr A(w)=0$ (second,
essential use of~(a)) and $\tr Q_2(w+k)=\tr Q_2(k)+O(|k|^3)$,
\[
\tr\big(M(v)-\ls I\big)=\tr Q_2(k)+O(|k|^3)+o(|k|^2)\le-\kappa|k|^2+o(|k|^2)<0,
\]
so $\lambda_{\min}M(v)\le\ls+\tfrac13\tr(M(v)-\ls I)<\ls$.
\end{proof}

The criterion is sufficient, not necessary; the sharper condition is
$\lambda_{\min}Q_2(k)<0$ on $K$. Note also that $\Lambda$ itself need not have a local maximum: off
$K$ it may increase at second order. Since $\Lambda=\tr M$, we have
$\tr Q_2(k)=\tfrac12D^2\Lambda(k,k)$.

\subsection{The Teichm\"uller chart}\label{subsec:chart}
For a word $w$ in the generators, let
\[
\tau_w(x):=\tr_{SO(2,1)}w\big(g_0(x),\dots,g_3(x)\big).
\]
We use the six words
\[
g_0g_1,\quad g_0G_2,\quad g_1g_2,\quad g_1g_3,\quad g_1G_3,\quad g_2g_3\qquad(G_k=g_k^{-1}).
\]

\begin{proposition}[Certified; Section~\ref{subsec:rank-cert}]\label{prop:G3}
The $6\times11$ Jacobian of $(\tau_w)_w$ on $T_{\xB}\Oct$ has rank $6$.
\end{proposition}

\begin{lemma}\label{lem:chart}
Let $\Sl\subset\Oct$ be the analytic $6$-dimensional coordinate slice through $\xB$ in the free
coordinates $\{3,6,10,12,13,15\}$ (Section~\ref{subsec:rank-cert}). Then:
\begin{enumerate}[label=\textup{(\alph*)}]
\item $\pi|_\Sl$ is a homeomorphism from a neighbourhood of $\xB$ in $\Sl$ onto an open neighbourhood
$N$ of $B$ in $\Teich$, and the functions $\bar\tau_w$ with $\tau_w=\bar\tau_w\circ\pi$ are local
coordinates of $\Teich$ at $B$.
\item For every word $w$ and every $C^1$ curve $x(t)$ in $\Oct$ through $\xB$,
$d(\tau_w\circ x)(v)=d\tau_w(\sigma(v))$, where $d\tau_w$ is the derivative of the
$\mathrm{Diff}_0$-invariant length-type function on smooth metrics. In particular $(d\tau_w)_w$ is
injective on TT tensors, and $\sigma\colon T_{\xB}\Sl\to T_B$ is an isomorphism.
\item There is an analytic retraction $p\colon\Oct\to\Sl$ near $\xB$ with $\pi\circ p=\pi$. Hence
$\Lambda_\Oct=\Lambda_\Sl\circ p$, $d\Lambda_\Sl(\xB)=0$,
$D^2\Lambda_\Oct(w,w)=D^2\Lambda_\Sl(dp\,w,dp\,w)$, and $\sigma\circ dp=\sigma$ on $T_{\xB}\Oct$.
\item Theorem~\ref{thm:F1} applies on $\Sl$: $M_\Sl(v)=\ls I+A(\sigma v)+Q_2(v)+O(|v|^3)$ and
$\tr(A\circ\sigma)\equiv0$.
\end{enumerate}
\end{lemma}

\begin{proof}
(a) $\tr_{SO(2,1)}$ is conjugation-invariant and needs no sign choice, so $\tau_w=\bar\tau_w\circ\pi$.
$\Phi:=(\tau_w)_w$ is analytic on $\Oct$ and its derivative on $T_{\xB}\Sl$ has rank $6$, so
$\Phi|_\Sl$ is an analytic diffeomorphism near $\xB$ onto an open $U\subset\R^6$. Hence $\pi|_\Sl$ is
continuous and injective. $\Teich$ is a topological $6$-manifold (Fenchel--Nielsen coordinates;
\cite[\S2]{FN}), so by invariance of domain $\pi|_\Sl$ is a homeomorphism onto an open $N\ni B$.

(b) With $k_t$ the smooth comparison family of Lemma~\ref{lem:N}, Step~4,
$\tau_w(x(t))=\tau_w(k_t)$, and $\dot k=\sigma(v)+L_Yh_0$. The function $\tau_w(k)$ is the trace of
the holonomy along a fixed loop of the $\mathfrak{so}(2,1)$-valued Cartan connection
$\omega_{LC}+\theta$ of $k$. This is defined and $C^1$ in $k$ for every metric, and it is flat, hence
invariant under free homotopy, exactly when the curvature is $-1$. So $\tau_w$ is a
$\mathrm{Diff}_0$-invariant $C^1$ function on hyperbolic metrics, and the Lie-derivative term
drops out. Rank $6$ then forces $\sigma$ to be onto the $6$-dimensional TT space, and injectivity
of $(d\tau_w)_w$ on TT follows.

(c) Set $p:=(\Phi|_\Sl)^{-1}\circ\Phi$. Since $\bar\tau$ is injective on $N$, $\pi\circ p=\pi$. As
$dp|_{T\Sl}=\mathrm{id}$ and $d\Lambda_\Oct(\xB)=0$, the Hessians correspond. Applying the injective
$(d\tau_w)_w$ to $d\tau(\sigma(dp\,v))=d(\tau\circ p)(v)=d\tau(\sigma(v))$ gives $\sigma\circ dp=\sigma$.

(d) $\Sl$ is an analytic sub-family of $\Oct$; use Lemma~\ref{lem:N}(a) and
Proposition~\ref{prop:trA}.
\end{proof}

\subsection{Invariance of the Hessian under all of $G$}

\begin{lemma}\label{lem:Ginv}
The quadratic form $H(\varsigma):=D^2\Lambda_\Sl(\sigma^{-1}\varsigma,\sigma^{-1}\varsigma)$ on TT tensors is
invariant under the natural action of $G$. Consequently $H|_{K_0}=c\cdot\gin$ for some $c\in\R$.
\end{lemma}

\begin{proof}
Let $g\in G$ act on $\Teich$ by change of marking, normalised by $g\cdot[k]=[(g^{-1})^*k]$. Then $g$
fixes $B$, and $\Lambda\circ g=\Lambda$ since the surfaces are isometric.

\emph{Step 1 ($g$ is analytic in the chart).} For every word $w'$,
$\bar\tau_{w'}(g\cdot X)=\bar\tau_{g^{-1}(w')}(X)$, because the $(g^{-1})^*k$-length of the class $w'$ is
the $k$-length of $g^{-1}(w')$. (For orientation-reversing $g$, $g^{-1}(w')$ is defined up to
conjugation and inversion, which does not affect $\tr_{SO(2,1)}$.) Since each $\tau_{w''}|_\Sl$ is
analytic and the chart is $(\Phi|_\Sl)^{-1}\circ\bar\tau$, the map $g_\Sl:=(\pi|_\Sl)^{-1}\circ g\circ\pi|_\Sl$ is
analytic near $\xB$ and fixes $\xB$. This uses continuity of the extended mapping class group
action on $\Teich$, so that $g\cdot\pi(x)\in N$ for $x$ near $\xB$. As $d\Lambda_\Sl(\xB)=0$, the Hessian is
invariant: $D^2\Lambda_\Sl(dg_\Sl v,dg_\Sl v)=D^2\Lambda_\Sl(v,v)$.

\emph{Step 2 ($\sigma$ intertwines).} Let $x(t)$ be a curve in $\Sl$ and $y=g_\Sl\circ x$. By
Lemma~\ref{lem:chart}(b), $d\tau_w(\sigma(y'))=\frac{d}{dt}\tau_{g^{-1}w}(x(t))=d\tau_{g^{-1}w}(\sigma(x'))$. On
hyperbolic metrics $\tau_{g^{-1}w}(k)=\tau_w((g^{-1})^*k)$; differentiating at the $g$-fixed $h_0$ gives
$d\tau_{g^{-1}w}(\varsigma)=d\tau_w(g\cdot\varsigma)$, since $(g^{-1})^*L_Yh_0=L_{g_*Y}h_0$ is killed by
$\mathrm{Diff}_0$-invariance. Injectivity of $(d\tau_w)_w$ on TT gives $\sigma(y')=g\cdot\sigma(x')$.

Steps~1 and~2 give $H(g\cdot\varsigma)=H(\varsigma)$ for all $g\in G$, in particular for $t$.
Proposition~\ref{prop:scalar} gives the last claim.
\end{proof}

\begin{remark}
An alternative route would cite real-analyticity of $\pi$ into the standard analytic structure of
$\Teich$ and analyticity of the action of the stabiliser of $B$ in the extended mapping class
group with derivative equal to the natural action on $\operatorname{Re}Q(B)$. The argument above
avoids these inputs by working entirely in the certified chart $\Sl$.
\end{remark}

\subsection{The Ky Fan bound along $\FB$}

\begin{lemma}[Ky Fan trace principle]\label{lem:KyFan}
For every $p$ and every $4$-dimensional $V\subset H^1(B)$,
$\Lambda(p)\le\tr_V(q_p\,|\,m_p)$, with equality when $V$ is the span of $1$ and the cluster
eigenfunctions.
\end{lemma}

\begin{proof}
Let $\mu_0\le\dots\le\mu_3$ be the eigenvalues of $q_p|_V$ relative to $m_p$. For each $j$, the span of
the first $j+1$ of their eigenvectors is a $(j+1)$-dimensional subspace on which the Rayleigh
quotient is at most $\mu_j$, so by min-max (\cite[Thm.~1.22]{Coo18}, from Chavel)
$\lambda_j\le\mu_j$. Summing and using $\lambda_0=0$ gives the claim. (For the finite-dimensional
statement see~\cite{Mos}, attributing it to~\cite{Fan49,Fan50}.)
\end{proof}

Take $V(p)=\operatorname{span}\{1,e_i+p\psi_i\}$, where $e_i\in E$ are exact eigenfunctions and
$\psi_i$ are fixed finite-element correctors, and set $g(p)=\tr(G_m(p)^{-1}G_q(p))$ with the Gram
matrices of $m_p,q_p$ on $V(p)$. Then $g\ge\Lambda\circ x$ and $g(0)=\Lambda(0)$. Both are analytic, so
$g-\Lambda\circ x\ge0$ has a minimum at $0$ and
\begin{equation}\label{eq:kyfan}
(\Lambda\circ x)''(0)\le g''(0).
\end{equation}
Even without twice differentiability, the $\limsup$ of the second difference quotients of
$\Lambda\circ x$ is at most $g''(0)$; the analyticity of Theorem~\ref{thm:F1} is needed only to read the
result as a statement about $D^2\Lambda$.

\begin{lemma}\label{lem:g2formula}
With $N=M_0^{-1}$,
\[
g''(0)=\tr(NX)-2\tr(NM_1NA)+2\ls\,c_1^{\mathsf T}Nc_1/m(1,1),
\]
where $X_{ij}=(q''-\ls m'')(e_i,e_j)+2b(\psi_i,e_j)+2b(e_i,\psi_j)+2(q-\ls m)(\psi_i,\psi_j)$,
$A_{ij}=b(e_i,e_j)$ with $b=q'-\ls m'$, $M_{1,ij}=m'(e_i,e_j)+m(\psi_i,e_j)+m(e_i,\psi_j)$,
$c_{1,i}=m'(e_i,1)+m(\psi_i,1)$ and $M_{0,ij}=m(e_i,e_j)$.
\end{lemma}

\begin{proof}
Expand $G_m^{-1}G_q$ to second order in $p$, using the Schur complement of the constant function,
$q_0(\psi,e)=\ls m_0(\psi,e)$, $Q(0)=\ls M_0$ and $m(e,1)=0$. The basis $e_i$ need not be orthonormal.
(The formula was re-derived independently and checked against an exact pencil test to relative
accuracy $3\times10^{-39}$, and against a five-point finite difference of $g$.)
\end{proof}

\begin{proposition}[Certified; Section~\ref{subsec:sign-cert}]\label{prop:sign}
Along $\FB$, $g''(0)\le-9.341154560446736$ ($n=96$, $P_2$ elements), with central value
$[-14.274\pm2.13\times10^{-4}]$. Consequently
\[
D^2\Lambda_\Oct(\vB,\vB)=(\Lambda\circ x)''(0)\le-9.3411<0.
\]
Under Hypothesis~\ref{hyp:SU} the same certificate gives $g''(0)\le-12.850476214196528$.
\end{proposition}

The last equality holds because $(\Lambda\circ x)''(0)=D^2\Lambda_\Oct(\vB,\vB)+d\Lambda_\Oct(x''(0))$ and
$d\Lambda_\Oct(\xB)=0$ (Proposition~\ref{prop:trA}).

%======================================================================
\section{Certified computations}\label{sec:cert}
%======================================================================

All computations below use Arb ball arithmetic through python-flint~0.9.0 (96- or 128-bit
precision), exact rationals or exact arithmetic in $\Q(\zeta_{24})$ (with sympy~1.14.0 for Sturm
sequences). Floating-point computations (numpy/scipy) only propose candidates. Every
truncation carries an explicit bound.

\subsection{Isolation of the cluster by the trace formula}\label{subsec:LP}

We use the Selberg trace formula in the normalisation of~\cite[(3.1)]{FBP-K}, which follows
Buser~\cite[p.~253]{Bus10}: for every admissible $f$ and closed hyperbolic surface $S$ of genus $g$,
\[
\sum_{\lambda\in\sigma(S)}\hat f(r(\lambda))=2(g-1)\int_0^\infty r\hat f(r)\tanh(\pi r)\,dr
+(2\pi)^{-1/2}\sum_{\gamma}\frac{\Lambda(\gamma)f(\ell(\gamma))}{2\sinh(\ell(\gamma)/2)},
\]
with the conventions of~\cite{FBP-K}. We use~\cite[Lemma~3.2]{FBP-K}: if $0<a\le\lambda_1(S)$, $a<b$,
$\hat f(r(\lambda))\ge c>0$ on $[a,b]$, $\hat f(r(\lambda))\ge0$ for $\lambda\ge b$, and $f(x)\le0$ for
$x\ge\operatorname{sys}(S)$, then the number of eigenvalues in $[a,b]$ is at most $c^{-1}(I+G_P-\hat
f(i/2))$, where $I$ is the integral term and $G_P$ the geodesic sum over any set $P$ of unoriented
primitive closed geodesics. We also use the following variant, whose proof is that
of~\cite[Lemma~3.2]{FBP-K} unchanged.

\begin{lemma}\label{lem:LemmaA}
Let $F(\lambda):=\hat f(r(\lambda))$. If $F\ge c>0$ on $[0,a]$, $F\ge0$ on $[a,\infty)$ and $f\le0$ on
$[\operatorname{sys}(S),\infty)$, then $\#(\sigma(S)\cap(0,a])\le(I+G_P-F(0))/c$.
\end{lemma}

\emph{Test functions.} $f(x)=u(x^2/s^2)e^{-x^2/(2s^2)}$ and $\hat f(\xi)=s\,v(s^2\xi^2)e^{-s^2\xi^2/2}$, with
$u(y)=\sum_{n\le N}c_nH_{2n}(\sqrt y)$, $v(y)=\sum_{n\le N}(-1)^nc_nH_{2n}(\sqrt y)$ ($H$ the physicists'
Hermite polynomials), $c_n,s\in\Q$ and $N=21$. Since $r(\lambda)^2=\lambda-\tfrac14$,
$F(\lambda)=s\,v(s^2(\lambda-\tfrac14))e^{-s^2(\lambda-1/4)/2}$. ($N$ must be odd: otherwise the leading
coefficients $c_N4^N$ and $(-1)^Nc_N4^N$ make $u\le0$ and $v\ge0$ near $+\infty$ incompatible.)

\emph{Verification.} The sign conditions $u\le0$ on $[y_0,\infty)$ (with rational $y_0\le\ell_1^2/s^2$) and
$v\ge0$ beyond the window are proved by exact Sturm counts: $0$ roots in each case. A lower bound $c$
for $F$ on the window comes from $4000$ Arb balls. The integral term uses exact Gaussian moments
$\int_0^\infty r(s^2r^2)^ke^{-s^2r^2/2}dr=2^kk!/s^2$ plus a rigorous \texttt{acb.integral} on $[0,12]$ with
an explicit tail bound. The geodesic term uses the lower bounds $12,12,24$ on the unoriented
multiplicities of $\ell_1=2\arccosh(1+\sqrt2)$ (the systole), $\ell_2=2\arccosh(3+2\sqrt2)$ and
$\ell_3=2\arccosh(5+3\sqrt2)$~\cite[\S7.4]{FBP-K}; these are primitive since $\ell_3<2\ell_1$, and only
lower bounds are needed because $f\le0$ there.

\begin{lemma}[Isolation and multiplicity]\label{lem:isolation}
\begin{enumerate}[label=\textup{(\alph*)}]
\item $\Spec(B)\cap(0,3.3]=\emptyset$ and $\#(\Spec(B)\cap[3.3,4.0])\le3$, with multiplicity.
\item $\lambda_3(B)\le3.8391848048$. Hence $\lambda_1,\lambda_2,\lambda_3\in(3.3,3.8391849]$ and
$\lambda_4>4.0$, a gap of at least $0.1608$. This part uses neither Hypothesis~\ref{hyp:SU} nor the
multiplicity theorem of~\cite{FBP-K}.
\item $\lambda_1=\lambda_2=\lambda_3\le3.8388872830$, $\dim E=3$, and $E$ is an irreducible real
$G$-module.
\item $\lambda_1\ge3.8388757010$ (Proposition~\ref{prop:LG}). Hence
$\lambda_1\in[3.8388757010,3.8388872830]$.
\end{enumerate}
\end{lemma}

\begin{proof}
(a) Run~A applies Lemma~\ref{lem:LemmaA} with $N=21$, $s=1$ on $[0,3.3]$: $c\ge0.9816424759937855$
and the certified bound is $[0.912612352675290375\pm9.95\times10^{-20}]<1$. Run~B
applies~\cite[Lemma~3.2]{FBP-K} with $a=3.3\le\lambda_1$ (from run~A), $N=21$, $s=13/10$ on
$[3.3,4.0]$: $c\ge0.9964283521820437$ and the bound is
$[3.86966111976973759\pm1.65\times10^{-18}]<4$.

(b) Section~\ref{subsec:FE}.

(c) By~\cite[Lemmas~7.2--7.3]{FBP-K} with $j=1$, whose proofs use symmetry, nodal domains
and~\cite[Prop.~4.4]{FBP-K} but no eigenvalue enclosure, the $\lambda_1$-eigenspace contains no
$1$- or $2$-dimensional irreducible real representation of $G$. So $m_1\ge3$, and with~(b),
$\lambda_1=\lambda_2=\lambda_3$ and $\dim E=3$; any decomposition of $E$ would contain a summand of
dimension $1$ or $2$, so $E$ is irreducible. The certified $P_2$
Rayleigh--Ritz bound $\lambda_3\le3.8388872829261946$ (Section~\ref{subsec:FE}) gives the upper end.

(d) Section~\ref{subsec:LG}; it uses only~(a).

\emph{Cross-check (not used).} By~(a) and~(b) the only eigenvalue in $(0,4.0]$ is $\lambda_1$, so the
eigenvalue of Hypothesis~\ref{hyp:SU}, if assumed, is $\lambda_1$; its interval is consistent with~(d).
\end{proof}

As sanity checks, windows that contain known eigenvalues give bounds above the true counts:
$3.4803\ge3$ for $[3.3,3.9]$ and $7.5466\ge7$ for $[3.3,5.4]$. The window end $X=4.05$ fails
(bound $4.0536$); with only three geodesic lengths the method does not reach beyond about $4.05$.

\subsection{Conforming finite elements on the glued octagon}\label{subsec:FE}

\emph{$P_1$ elements.} On each geodesic triangle of the mesh, trial functions are affine in the
barycentric parameters $(a,b)$ of Section~\ref{subsec:mesh}. Isometries act linearly on $W$, so the
construction is chart-free, and the trace on an edge depends only on the two endpoint values. The
space is therefore \emph{exactly conforming on the glued surface}, including across the projective
side pairings, it lies in $H^1(B)$, and it contains the constants. (Klein-coordinate $P_1$ elements
lack this property.)

\emph{$P_2$ elements.} Polynomials of degree $\le2$ in $(a,b)$. On an edge, $W=sP_i+(1-s)P_j$ and
the trace is a quadratic in $s$ determined by the two endpoint values and the value at $s=\frac12$,
the geodesic midpoint. Isometries preserve $s$, so this space is also exactly conforming and
contains the constants. Edge degrees of freedom are keyed by glued endpoint labels; the code
asserts Euler characteristic $-2$ and exactly two triangles per edge.

\emph{Element integrals.} All entries are integrals $\int_{\text{simplex}}p\cdot Q^{-s}$ with
$s\in\{\tfrac12,\tfrac32\}$ ($P_1$) or $s\in\{\tfrac12,\tfrac32,\tfrac52,\tfrac72\}$ ($P_2$). Write
$Q=Q_c(1+\epsilon)$ and expand $(1+\epsilon)^{-s}$ in a truncated binomial series ($4$ terms for $P_1$, $6$
for $P_2$) with a Lagrange remainder, e.g.\ for $P_1$
$|R|\le|\binom{-s}{4}|\epsilon_{\max}^4(1-\epsilon_{\max})^{-s-4}$; monomials are integrated exactly by
$\int a^ib^j=i!\,j!/(i+j+2)!$.

\emph{Rayleigh--Ritz.} Untrusted float eigenvectors $x_1,x_2,x_3$ and the constant span a
$4$-dimensional $V$; its $4\times4$ matrices $K,M$ are summed exactly in Arb. By min-max
\cite[Thm.~1.22]{Coo18}, $\lambda_3(B)\le\mu_3(K,M)$, and $\mu_3<t$ is certified by Sylvester's
criterion for $tM-K$, with $t$ found by bisection. Applied to the same space, Lemma~\ref{lem:KyFan} gives a certified upper bound for
$\lambda_1+\lambda_2+\lambda_3$.

The certified bounds are in Tables~\ref{tab:P1} and~\ref{tab:P2}. At $n=192$ ($P_1$),
$\lambda_3(B)\le3.8391848048$; at $n=96$ ($P_2$), $\lambda_3(B)\le3.8388872830$ and
$\lambda_1+\lambda_2+\lambda_3\le11.516661838577997$. The float $P_2$ errors decrease like $h^4$
($4.8\times10^{-6}$, $3.0\times10^{-7}$, $1.9\times10^{-8}$). The discrete cluster splits as $2+1$
because the mesh has only the octagon's symmetry.

\subsection{A guaranteed lower bound for \texorpdfstring{$\lambda_1$}{lambda1}: Lehmann--Goerisch}\label{subsec:LG}

Rayleigh--Ritz gives only upper bounds. For the lower bound we use the Lehmann--Goerisch method in
the form of~\cite[\S4.1, Thms.~4.1 and~4.3]{LiuPlum}, which goes back to Goerisch; the same statement,
written with an explicit shift, is~\cite[Thm.~2.1]{BMPW} as reproduced in~\cite[\S2]{Vej18}. We quote it
for our setting. Let $D$ be a real vector space with symmetric bilinear forms $M$ (positive definite)
and $N$ whose eigenvalue problem $M(u,\varphi)=\tilde\lambda N(u,\varphi)$ has a discrete spectrum
$\tilde\lambda_1\le\tilde\lambda_2\le\cdots$ with an $M$-orthonormal eigenbasis. Let $X$ be a real vector space
with a positive semidefinite symmetric form $b_G$ and $T\colon D\to X$ linear with $b_G(Tu,Tv)=M(u,v)$. Given
$v_1,\dots,v_n\in D$, let $w_1,\dots,w_n\in X$ satisfy $b_G(w_i,T\varphi)=N(v_i,\varphi)$ for all $\varphi\in D$.
For $\rho>0$ put $A_0=(M(v_i,v_j))$, $A_1=(N(v_i,v_j))$, $A_2=(b_G(w_i,w_j))$, $A=A_0-\rho A_1$ and
$B=A_0-2\rho A_1+\rho^2A_2$, and suppose $B$ is positive definite; let $\nu_1\le\dots\le\nu_n$ be the
eigenvalues of $Az=\nu Bz$ and $q$ the number of negative ones.

\begin{theorem}[Lehmann--Goerisch; quoted from~\cite{LiuPlum}]\label{thm:LG}
If $\rho\le\tilde\lambda_{m+1}$ for some $m\ge n$, then
$\tilde\lambda_{m+1-k}\ge\rho-\rho/(1-\nu_k)$ for $1\le k\le q$.
\end{theorem}

\emph{Instantiation.} Take $D=H^1(B)$, $M=q+m$ and $N=m$. The eigenvalues are shifted,
$\tilde\lambda=\lambda+1$: $\tilde\lambda_1=1$ (the constants), $\tilde\lambda_{2,3,4}=\lambda_{1,2,3}+1$ and
$\tilde\lambda_5=\lambda_4+1$. By Lemma~\ref{lem:isolation}(a), $\lambda_4>4.0$, so $\rho:=5\le\tilde\lambda_5$; we take
$m=4$ and $n=3$. Nothing else about the spectrum is used. In particular neither
Hypothesis~\ref{hyp:SU} nor~\cite[Lemmas~7.2--7.3]{FBP-K} enters. Let $X=L^2(TB)\times L^2(B)$,
$Tu=(\nabla u,u)$, and let $b_G$ be the $L^2$ product, so $b_G(Tu,Tv)=M(u,v)$. For trial functions $v_i$ and
any $\sigma_i\in H(\mathrm{div},B)$ put $w_i=(\sigma_i,\,v_i+\operatorname{div}\sigma_i)$. Integration by parts on the
closed surface gives, for all $\varphi\in H^1(B)$,
\[
b_G(w_i,T\varphi)=\int\sigma_i\cdot\nabla\varphi+\int(v_i+\operatorname{div}\sigma_i)\varphi=\int v_i\varphi=N(v_i,\varphi),
\]
exactly and with no constraint on $\operatorname{div}\sigma_i$. This is why the shift is used. The choice of
$\sigma_i$ affects only the sharpness. With $k=3$, $m+1-k=2$, so the theorem bounds
$\tilde\lambda_2=\lambda_1+1$, the first \emph{nonzero} eigenvalue. Since $\lambda_2,\lambda_3\ge\lambda_1$, the bound
holds for all three.

\emph{Discretisation.} The $v_i$ are conforming $P_2$ functions on the $n=48$ mesh of
Section~\ref{subsec:FE} (untrusted float eigenvectors, used as exact data). The $\sigma_i$ lie in a
Raviart--Thomas space of degree $2$ in each reference simplex, with unknown $j=\sqrt g\,\sigma$.
Its degrees of freedom are the outward flux densities per unit edge parameter $s\in\{0,\frac12,1\}$
and six interior moments. The flux through an edge segment is the intrinsic $1$-form
$\iota_\sigma dA$, and the edge point $W=sP_i+(1-s)P_j$ is the same from both sides, also across the
projective side pairings, since isometries are linear on $W$ and preserve $s$. Equal and opposite
edge degrees of freedom therefore make $\sigma_i$ exactly $H(\mathrm{div})$-conforming on the glued
surface, by the same mechanism as $P_2$ conformity. In the reference chart,
$|\sigma|^2dA=j^{\mathsf T}(LL^{\mathsf T}-QG)j/(D\sqrt Q)\,da\,db$ and $(\operatorname{div}\sigma)\sqrt g=\partial_aj^a+\partial_bj^b$
with $\sqrt g=DQ^{-3/2}$. The entries of $A_0,A_1,A_2$ are sums over the $18\,432$ elements of integrals
$\int p\cdot Q^{-s}$ ($s\in\{\frac12,\frac32\}$), enclosed as in Section~\ref{subsec:FE} with $14$ series terms
and a Lagrange remainder. The $\sigma_i$ are float minimisers of $b_G(w_i,w_i)$ (a sparse
positive definite solve).

\begin{proposition}[Certified]\label{prop:LG}
In Arb, $B\succ0$ and $\bar\nu B-A\succ0$ with $\bar\nu=-30.031942622361008<0$, by Sylvester's
criterion. Hence $\nu_1,\nu_2,\nu_3<\bar\nu<0$, $q=3$, and Theorem~\ref{thm:LG} with $k=3$ gives
\[
\lambda_1(B)\ge4-\frac{5}{1-\bar\nu}\in[3.8388757010527243542306205221870856407\pm6.06\times10^{-38}],
\]
so $\lambda_1(B)\ge3.8388757010527243$ (rounded down).
\end{proposition}

At $n=24$ the same construction gives only $3.8387028952$. The float $\nu$ at $n=48$ are
$-30.03241$, $-30.03241$ and $-30.03194$.

\emph{Checks.} (i) On the flat torus $\R^2/\Z^2$ ($\lambda_1=4\pi^2$ with multiplicity $4$; $n=4$ trial
functions, $\rho=79.9\le1+8\pi^2$) the same code certifies $39.2041077545$, $39.4582495177$,
$39.4770992124$ and $39.4783342462$ on $N=4,8,16,32$ meshes, all below $4\pi^2=39.4784176044$ and
converging at rate $O(h^4)$. (ii) The global Green identity
$\sum_T\int(\sigma\cdot\nabla v+(\operatorname{div}\sigma)v)=0$, evaluated in Arb for the certified fluxes,
has residual at most $7.4\times10^{-36}$. (iii) Making the fluxes non-conforming (all edge signs $+1$) makes
$B$ indefinite on the torus, as the theory predicts. (iv) An independent reviewer redid the final step
from the stored enclosures of $A_0,A_1,A_2$ in mpmath interval arithmetic and obtained the same bound. The
reviewer also reproduced the element integrals from a from-scratch metric to relative accuracy
$1.2\times10^{-32}$, checked every glued edge geometrically, and confirmed on an exact toy model that wrong
indexing of the zero mode, or an invalid $\rho$, produces false bounds.

\emph{The hypothesis on $\rho$ is not self-checking.} With $\rho=6.5$, or $\rho=7.5$ at $n=24$, both
above $\lambda_4+1$, the computation still returns values below the true $\lambda_1$. A wrong $\rho$ would
therefore go undetected. The trace-formula certificate of Lemma~\ref{lem:isolation}(a) is load-bearing.

\emph{Scope.} The interval $[3.8388757010,3.8388872830]$ does not by itself imply multiplicity~$3$, so
Lemma~\ref{lem:isolation}(c) and the certificates below still use~\cite[Lemmas~7.2--7.3]{FBP-K}.

\subsection{From finite-element functions to exact eigenfunctions}\label{subsec:transfer}

\begin{lemma}[Gap transfer]\label{lem:transfer}
Let $u\in H^1(B)$ with $m(u,1)=0$, let $e=P_Eu$ and $w=u-e$, and let $\mu=q(u)/m(u)$. If
$\lambda_{\mathrm{low}}\le\lambda_1$ and $\lambda_4\ge\Lambda_4>\lambda_1$, then
\[
m(w)\le\frac{\mu-\lambda_{\mathrm{low}}}{\Lambda_4-\lambda_{\mathrm{low}}}\,m(u),\qquad
q(w)\le\frac{(\mu-\lambda_{\mathrm{low}})\Lambda_4}{\Lambda_4-\lambda_{\mathrm{low}}}\,m(u).
\]
\end{lemma}

\begin{proof}
$w$ is orthogonal to $E\oplus\R1$ in both $m$ and $q$, so $q(w)\ge\lambda_4m(w)$ and
$q(u)-\lambda_1m(u)=q(w)-\lambda_1m(w)\ge\max\{(\lambda_4-\lambda_1)m(w),(1-\lambda_1/\lambda_4)q(w)\}$.
Since $\mu<\Lambda_4$, both right-hand sides decrease when $\lambda_{\mathrm{low}}$ is replaced by
$\lambda_1$ or $\Lambda_4$ by $\lambda_4$, so the stated bounds follow.
\end{proof}

This is where the lower bound for $\lambda_1$ enters the certificates of
Sections~\ref{subsec:Anonzero-cert} and~\ref{subsec:sign-cert}; we use $\Lambda_4=4.0$ and
$\lambda_1$ in the ball
\[
3.83888149202636215\pm5.791\times10^{-6}=[3.838875701026362,\,3.838887283026362].
\] This ball contains
$[3.8388757010527243,\,3.8388872829261946]$, the interval given by Proposition~\ref{prop:LG} and the $P_2$
bound. The scripts differ from the original ones only in this ball and the output file name.

\emph{Alternative with Hypothesis~\ref{hyp:SU}.} The original runs used the ball
$3.8388872588\pm1.0000001\times10^{-6}$. That ball's lower
end lies $4.2\times10^{-11}$ below the lower end of Hypothesis~\ref{hyp:SU} (conservative), and its
upper end, $4.2\times10^{-11}$ below that of Hypothesis~\ref{hyp:SU}, is covered by the certified
$P_2$ bound $\lambda_1\le3.8388872830$. A re-run of the sign certificate with the exact lower end
$3.83888625884$ gives $-12.850504$ ($n=96$).

\subsection{The certificate for $A\ne0$}\label{subsec:Anonzero-cert}
The pulled-back forms along the deformation $x(t)$ of Proposition~\ref{prop:Anonzero} are those
of~\eqref{eq:forms}, and $b=q'-\lambda_1m'$ is computed elementwise by forward-mode dual numbers over
Arb, using $\frac{d}{dt}\int pQ^{-s}=\int\dot pQ^{-s}-s\int p\dot QQ^{-s-1}$. Let $x$ be the (untrusted)
float $P_1$ eigencombination minimising $b$, and $u=x-(m(x,1)/m(1,1))\cdot1$, exactly
$m$-orthogonal to constants. With $e=P_Eu$ and $w=u-e$,
\[
b(e,e)=b(u,u)-2b(u,w)+b(w,w),
\]
and
$|b(u,w)|\le(\sum_T\beta_T^2q_T(u))^{1/2}q(w)^{1/2}+\lambda_1(\sum_T\beta_{\rho,T}^2m_T(u))^{1/2}m(w)^{1/2}$,
$|b(w,w)|\le\beta_{\max}q(w)+\lambda_1\beta_{\rho,\max}m(w)$, where $-\beta_Ta_T\le\dot a_T\le\beta_Ta_T$ and
$|\dot\rho_T|\le\beta_{\rho,T}\rho_T$ are enclosed in Arb over a parameter box containing the simplex.
Here $\beta_T\ge\|N^{-1}\dot N\|_F+|\dot D/D+\dot Q/(2Q)|$ and $\beta_{\rho,T}\ge|\dot D/D-1.5\,\dot Q/Q|$.
Finally $m(e)=m(u)-m(w)>0$, so $e\ne0$. The results are in Table~\ref{tab:Anonzero}; at $n=192$,
$b(e,e)\le-3.0383476242423058\cdot m(u)$ with the ball of Section~\ref{subsec:transfer}
($-3.0875694043934345\cdot m(u)$ with the ball of Hypothesis~\ref{hyp:SU}). An independent re-run
without the symmetry shortcut gives $-3.0383476316928864\cdot m(u)$.

The printed $\beta_{\max}\approx33$ is mostly enclosure overestimation (Frobenius norm and
dependency over a box containing the simplex), not geometry: at the side-midpoint elements the
sampled true supremum is about $3.3$. Overestimation is safe. The certificate survives any
$\lambda_1$ lower bound $\ge3.83788$; with $\lambda_4\ge3.9$ it still gives $-1.89$; and a mesh-only
deformation, for which $A=0$ exactly, is correctly not certified ($+1.409$).

\emph{Symmetry shortcut.} The deformation and the mesh are invariant under the rotation by
$\pi/2$, and the eight sectors are congruent blocks with identical local node order, so only two
sectors are computed. This is justified by exact invariance and was checked against the full
computation at $n=48$ (identical output) and by an independent computation of all $4608$ elements
at $n=16$ ($55\,296$ pairs of entry enclosures, none non-overlapping).

\subsection{The sign certificate}\label{subsec:sign-cert}

\emph{Correctors (untrusted).} $\psi_i$ solves the discrete first-order system
$(K-\bar\lambda M)\psi_i=-(K'-\bar\lambda M')x_i$ on the $M$-complement of $\{x_1,x_2,x_3,1\}$, then
adds $-\tfrac12\sum_jm'(x_i,x_j)x_j$ (so $M_1\approx0$) and $-m'(x_i,1)/m(1,1)$ (so $c_1\approx0$). These
choices only sharpen the bound.

\emph{Jets.} Everything is computed with order-$2$ Taylor jets in $p$ over Arb: the vertices
($\FB$ in closed form), the mesh nodes, $N=\operatorname{adj}(LL^{\mathsf T}-QG)$, $D$, $Q$ and the
integrals, using
$Q(p)^{-s}=Q_0^{-s}[1-s(Q_1/Q_0)p+(-sQ_2/Q_0+\tfrac{s(s+1)}2(Q_1/Q_0)^2)p^2+\cdots]$.
Element derivatives agree with central differences to $O(h^2)$.

\emph{Transfer.} Put $u_i=x_i-c_i1$ with $c_i=m(x_i,1)/m(1,1)$ enclosed in Arb, $e_i=P_Eu_i$ and
$w_i=u_i-e_i$. Lemma~\ref{lem:transfer} bounds $w_i$ (Table~\ref{tab:eta}). Every entry of $X,A,M_1,c_1,M_0$
at $e$ is its value at $u$ plus cross terms in $w$, bounded using elementwise relative bounds
$-\beta^1a\le\dot a\le\beta^1a$, $-\beta^2a\le\ddot a\le\beta^2a$ (and similarly for $\rho$), enclosed over
three boxes covering the simplex; quadratic terms in $w$ use $\beta_{\max}$. The $3\times3$ algebra is
done in Arb.

\emph{Result.} At $n=96$, with the ball of Section~\ref{subsec:transfer}:
$g''(0)\le-9.341154560446736$, central value $[-14.274\pm2.13\times10^{-4}]$ (an independent re-run gives
$-9.341154566965994$). The distances to $E\oplus\R1$ are larger than with Hypothesis~\ref{hyp:SU}
(Table~\ref{tab:eta}), because they are dominated by the width of the $\lambda_1$ ball.

\emph{Alternative with Hypothesis~\ref{hyp:SU}.} At $n=96$:
$g''(0)\le-12.850476214196528$, central value $[-14.2742\pm2.23\times10^{-5}]$; the joint transfer
error is about $1.42$, so the margin is about nine times the error. The rest of this subsection
refers to these runs. Table~\ref{tab:sign} lists the
results for $n=24,48,96$ and Table~\ref{tab:budget} the error budget. The breakdown rows of
Table~\ref{tab:budget} are ``one at a time'' (each error group switched on alone); they do not add
up to the joint error, and adding them to the central value would give about $-12.864$, which is
\emph{below} the certified value. Only the joint ball is a bound.

\emph{Robustness.} The symmetry shortcut is not needed: runs without it certify at $n=24$ and
$n=48$ ($-12.566084126010532$), and an independent re-run without it at $n=96$ gives
$-12.850476213730866$. The $\beta$ values are enclosure values: on about $150$ sampled elements the
ratio of the observed to the enclosed value is at most $0.70$ ($\beta^1$), $0.89$ ($\beta^2$), $0.99$
($\beta^1_\rho$) and $0.97$ ($\beta^2_\rho$).

\emph{Sensitivity to the $\lambda_1$ lower bound.} The $n=96$ certificate holds for every rigorous
lower bound $\lambda_1\ge3.8388029653$ and fails at $3.8388029297$. The bound of
Proposition~\ref{prop:LG} clears this threshold by about $7.3\times10^{-5}$.

\emph{Negative controls.} Each of the following fails to certify, as it should: $\psi=0$;
$\lambda$-radius $10^{-4}$; $\lambda_4=3.84$; a $\lambda$-ball excluding the true eigenvalue.

\emph{Independent floating-point checks (not part of the proof).} A separately written conforming
$P_1$ code on a different mesh, with re-meshed finite differences, gives $\Lambda''\approx-14.26$ by
Richardson extrapolation at $n=64/128$; the five-point difference of the discrete $P_2$ cluster sum
gives $-14.2739$ at $n=24$.

\subsection{Rank certificates}\label{subsec:rank-cert}

\begin{proposition}\label{prop:rankJ}
The $5\times16$ constraint Jacobian $J$ at $\xB$ has rank $5$: its minor in columns $(0,2,5,7,8)$ has
determinant $[46151927.03406868578682884302258\pm2.89\times10^{-24}]$. For the regular octagon,
$d(\text{angle sum})/dR=[-5.572971195984104378981991909236\pm5.18\times10^{-31}]$, in agreement with
the closed form $8\,d\theta/dR$.
\end{proposition}

\emph{Proof of Proposition~\ref{prop:G3}.} The pairings $g_k$ are computed in Arb at $128$ bits
from the frame formula of Section~\ref{subsec:octspace}, with derivatives by forward-mode dual
numbers (no finite differences). At $\xB$, $|g_kV_{k+4}-V_{k+1}|,|g_kV_{k+5}-V_k|\le1.2\times10^{-29}$,
$|\mathsf B g_k^{\mathsf T}\mathsf Bg_k-I|\le4.4\times10^{-29}$, the constraints enclose $0$, and
$\tr g_k$ overlaps $1+2\cosh(2\arccosh(1+\sqrt2))=22.3137084989848$, so each $g_k$ translates by the
systole. The six traces at $\xB$ are $340.705627484771$ ($g_0g_1,g_1g_2,g_2g_3$) and
$134.882250993909$ ($g_0G_2,g_1g_3,g_1G_3$). A basis of $\ker J$ is formed exactly in Arb from the
pivot minor of Proposition~\ref{prop:rankJ}, $T_b=[-J_{\mathrm{piv}}^{-1}J_{\mathrm{free}};I]$, using the
remaining $11$ coordinates as free parameters. The resulting $6\times11$ Jacobian has $392$ of its
$462$ maximal minors certified nonzero; the best uses tangent columns $(1,3,5,7,8,10)$, i.e.\ the
free coordinates $3,6,10,12,13,15$, with determinant $[7872428136359854080.0\pm3\times10^{-6}]$. As a
cross-check, an $11\times11$ minor of the stacked $[J;J_f]$ has determinant
$[-1.3827971011809151404\times10^{24}\pm2.74\times10^4]$. \qed

The four generator traces $\tau_{g_k}$ have rank only $3$ on $T_{\xB}\Oct$, and several
symmetric-looking six-word sets have rank $5$ at $B$; these coincidences are forced by the octagon
symmetry, which is why the specific words above are used. In an orthonormal tangent basis the six
rows have singular values between $5.7\times10^3$ and $2.65\times10^2$ (floating point).

%======================================================================
\section{Proof of the Main Theorem}\label{sec:proof}
%======================================================================

\begin{proof}
Item~(i) is Lemma~\ref{lem:isolation}, and item~(ii) is Corollary~\ref{cor:rank} with
Section~\ref{sec:symmetry}. Work in the chart $\Sl$ of Lemma~\ref{lem:chart}.
\begin{enumerate}
\item $\Lambda_\Sl$ is real-analytic, and $M_\Sl$ has the expansion of Theorem~\ref{thm:F1} with first-order
term $A\circ\sigma$ (Lemma~\ref{lem:N}(a), Lemma~\ref{lem:chart}(d)); $\tr(A\circ\sigma)\equiv0$
(Proposition~\ref{prop:trA}).
\item $K_\Sl:=\ker(A\circ\sigma)=\sigma^{-1}(K_0)$ is $3$-dimensional (Corollary~\ref{cor:rank};
$\sigma$ is an isomorphism by Lemma~\ref{lem:chart}(b)).
\item $dp\,\vB\in K_\Sl\setminus\{0\}$, since $\sigma(dp\,\vB)=\sigma(\vB)=\wB\in K_0\setminus\{0\}$
(Lemmas~\ref{lem:FBline} and~\ref{lem:chart}(c)), and
$D^2\Lambda_\Sl(dp\,\vB,dp\,\vB)=D^2\Lambda_\Oct(\vB,\vB)\le-9.3411$ (Proposition~\ref{prop:sign},
Lemma~\ref{lem:chart}(c)). This is item~(iv).
\item By Lemma~\ref{lem:Ginv}, $D^2\Lambda_\Sl|_{K_\Sl}=c\cdot\sigma^*\gin$. By step~3,
$c\,\gin(\wB,\wB)\le-9.3411$, so $c<0$ and $D^2\Lambda_\Sl|_{K_\Sl}$ is negative definite. This is
item~(iii), and hypothesis~(b) of Lemma~\ref{lem:criterion} holds with $\kappa>0$.
\item Lemma~\ref{lem:criterion} gives $\lambda_1(\pi(x))=\lambda_{\min}M_\Sl(x)<\ls$ for $x\in\Sl$ near $\xB$,
$x\ne\xB$. Case~(i) of its proof gives the first-order statement of item~(iii).
\item Since $\pi|_\Sl$ is a homeomorphism onto a neighbourhood $N$ of $B$ (Lemma~\ref{lem:chart}(a)),
$\lambda_1<\lambda_1(B)$ on $N\setminus\{B\}$.\qedhere
\end{enumerate}
\end{proof}

\begin{table}[!htbp]
\caption{External inputs. ``Quoted'' means a published statement used without re-proof.}
\label{tab:deps}
\small
\begin{tabular}{@{}p{0.5\textwidth}p{0.18\textwidth}p{0.24\textwidth}@{}}
\toprule
Input & Used in & Status\\
\midrule
Lehmann--Goerisch theorem~\cite[Thms.~4.1, 4.3]{LiuPlum}; second source~\cite[Thm.~2.1]{BMPW} as reproduced in~\cite[\S2]{Vej18} &
Proposition~\ref{prop:LG}, hence Lemma~\ref{lem:transfer} & quoted from secondary sources; Goerisch's original papers not read\\
Selberg trace formula in the normalisation of~\cite[(3.1)]{FBP-K} (from~\cite{Bus10});
\cite[Lemma~3.2]{FBP-K} and its proof & Lemma~\ref{lem:isolation}(a) & quoted\\
Systole $2\arccosh(1+\sqrt2)$ and multiplicities $\ge12,12,24$~\cite[\S7.4]{FBP-K} & Lemma~\ref{lem:isolation}(a), hence $\rho=5$ in Proposition~\ref{prop:LG} & quoted; lower bounds only\\
No $1$- or $2$-dimensional irreducibles in the $\lambda_1$-eigenspace~\cite[Lemmas~7.2--7.3]{FBP-K} & Lemma~\ref{lem:isolation}(c) (not Proposition~\ref{prop:LG}) & quoted\\
$|\Isom(B)|=96$, $\Isom^+(B)\cong GL(2,3)$, octahedral quotient~\cite[\S7.1]{FBP-K}; finite $S_4\subset PSL(2,\C)$ is octahedral & Section~\ref{sec:symmetry} & quoted / standard\\
$T_B\Teich\cong\{\operatorname{Re}\varphi\}$~\cite{Maz,Tro92,DW} & Lemma~\ref{lem:N}, Section~\ref{sec:symmetry} & quoted\\
Fischer--Marsden splitting (as stated in~\cite[Rem.~3.3]{Maz}) & Lemma~\ref{lem:N} & quoted\\
Min-max (\cite[Thm.~1.22]{Coo18}, from Chavel); Ky Fan~\cite{Mos,Fan49,Fan50} & Lemmas~\ref{lem:isolation}, \ref{lem:KyFan} & quoted; Ky Fan reduced to min-max\\
Elliptic regularity, Rellich compactness, spectral theorem, Neumann series, Riemannian centre of
mass in $\Hyp$, Poincar\'e polygon theorem, invariance of domain & Sections~\ref{sec:setup}--\ref{sec:second} & standard\\
$\Teich$ is a topological $6$-manifold~\cite{FN}; continuity of the extended mapping class group action & Lemmas~\ref{lem:chart}, \ref{lem:Ginv} & quoted / standard\\
python-flint~0.9.0 (Arb), sympy~1.14.0 & Section~\ref{sec:cert} & trusted software\\
\midrule
\emph{Optional, not used:} Hypothesis~\ref{hyp:SU}~\cite[\S5]{SU13} & cross-check; alternative constants & quoted\\
\bottomrule
\end{tabular}
\end{table}

%======================================================================
\section{Formal verification in Lean}\label{sec:lean}
%======================================================================

We formalised the finite-dimensional part of the argument in Lean~4 (toolchain v4.34.1, with
Mathlib). The development has about $2250$ lines in $8$ files and contains no \texttt{sorry}, no
\texttt{axiom} declarations and no \texttt{native\_decide}. Its main theorem is
\begin{align*}
&\texttt{theorem Bolza.bolza\_strict\_local\_max (I : BolzaInputs)}:\\
&\qquad\forall^{\mathrm{f}}X\in\mathcal N(\texttt{I.B}),\ \ X\ne\texttt{I.B}\ \Rightarrow\ \texttt{I.lam1}\,X<\texttt{I.lam1}\,\texttt{I.B},
\end{align*}
where $\forall^{\mathrm f}X\in\mathcal N(\cdot)$ is Mathlib's ``eventually in the neighbourhood filter''
and \texttt{\#print axioms} reports only \texttt{propext}, \texttt{Classical.choice} and
\texttt{Quot.sound}.

\subsection{Scope}
Every analytic, geometric and computational input is a named field (an explicit hypothesis) of the
structure \texttt{BolzaInputs}. These fields encode: $\Teich$ as a topological space and the chart
of Lemma~\ref{lem:chart}(a); the cluster matrix of Theorem~\ref{thm:F1} ($C^2$, symmetric, $M(0)=\ls I$,
$\lambda_1\le$ Rayleigh quotients); the representation of $G$ on $E$ and the equivariance of
Lemma~\ref{lem:N}(b); irreducibility of $E$; the group relations used in
Proposition~\ref{prop:K0kerA}; $A\ne0$; the invariance of Lemma~\ref{lem:Ginv}; the vector $\wB$ and
its symmetries; the curve $\FB$ in the chart; the operator Ky Fan comparison; and the certified
value of $g''(0)$.

Machine-checked from these inputs are: the local-maximum criterion (Lemma~\ref{lem:criterion}) in
$C^2$ form, including a uniform second-order Taylor theorem; $\tr A\equiv0$; $K_0\subset\ker A$,
$\ker A\subset K_0$ and $D^2\Lambda|_{K_0}=c\,I$; the exact line of Lemma~\ref{lem:FBline}; the step
$c<0$; the finite-dimensional Ky Fan theorem; the step $(\Lambda\circ\gamma)''\le g''$ from the Ky Fan
comparison; and the second-order chain rule turning it into $D^2\Lambda(\wB,\wB)\le-12.8504$. The exact
representation theory (the $48$-element image group, characters and multiplicities, generator
matrices, change of basis) is checked by kernel evaluation. The final arithmetic of each certificate
(e.g.\ $0.912612<1$, $3.869661<4$, $-12.850476214196528\le-12.8504$, the determinant enclosures
excluding $0$) is checked by \texttt{norm\_num}. A toy instance of \texttt{BolzaInputs} shows that
the hypotheses are jointly satisfiable; it says nothing about $B$.

\subsection{Limits}
The Lean development does \emph{not} verify: that the fields of \texttt{BolzaInputs} hold for the
Bolza surface (this is the content of Theorem~\ref{thm:F1}, Lemmas~\ref{lem:N}, \ref{lem:chart}
and~\ref{lem:Ginv}, the certificates, and the quoted results); the Arb computations themselves;
the operator min-max and Ky Fan inequality on $H^1$; or the identification of the exact matrices
with the geometric action on TT tensors. The step from character norm $1$ to absolute
irreducibility is not formalised; instead, the facts actually used (Schur-type invariant forms on
$K_0$, the absence of invariant subspaces in $W_0$, and the involution argument) are proved
directly over $\R$.

The Lean development was written against the constant $-12.8504$ of the run under
Hypothesis~\ref{hyp:SU}. The argument uses that constant only through $c<0$, so the SU-free constant
$-9.3411$ plays the same role. The corresponding \texttt{norm\_num} check
($-9.341154560446736\le-9.3411$) and the Lehmann--Goerisch step have not been added to the
formalisation.

\subsection{Observations from the formalisation}
\begin{enumerate}
\item The hypothesis $A\ne0$ is not needed for Lemma~\ref{lem:criterion}; symmetry of $M$ is used only
to pass to an eigenvalue statement. The constant $c_A$ comes from an elementary bound: a traceless
symmetric $3\times3$ matrix $S$ has a test vector with entries in $\{0,\pm1\}$ and $x^{\mathsf T}Sx\le-\|S\|_{\sup}/4$.
\item $\tr A\equiv0$ needs equivariance under $r$ only.
\item Both $\ker A=K_0$ and the scalar form of the Hessian genuinely use the non-octagon element
$t$ (Section~\ref{subsec:D8}).
\item As noted after Table~\ref{tab:budget}, the one-at-a-time error rows of the sign certificate
are not additive; only the joint enclosure is checked.
\end{enumerate}

%======================================================================
\section{Remarks and open problems}\label{sec:remarks}
%======================================================================

\begin{remark}[The input from~\cite{SU13}, and the weakest inputs]\label{rem:SU}
An earlier version of the proof depended on Hypothesis~\ref{hyp:SU}. It rests on a single sentence
of~\cite[\S5]{SU13} reporting a rigorous implementation whose code is unpublished, and the
sentence does not say whether the error $10^{-6}$ is absolute or relative. Proposition~\ref{prop:LG}
removes this dependence. Hypothesis~\ref{hyp:SU} remains consistent with our enclosure
$[3.8388757010,3.8388872830]$. If it is assumed, the constants improve to $-12.8504$ (second order)
and $-3.0875\cdot m(u)$ (first order). The second-order certificate needs only
$\lambda_1(B)\ge3.8388029653$; under Hypothesis~\ref{hyp:SU} this leaves $83$ times its stated error,
and it still holds if that error is read as relative ($n=48$, radius $10^{-5}$, gives $-9.43$).

The weakest inputs are now quoted results rather than computations:~\cite[Lemmas~7.2--7.3]{FBP-K}
(for multiplicity $3$), the geodesic data of~\cite[\S7.4]{FBP-K} with the trace formula (for
$\lambda_4>4.0$, which is the load-bearing hypothesis $\rho=5$ of Proposition~\ref{prop:LG}), and the
Lehmann--Goerisch theorem. We quote the latter from secondary sources~\cite{LiuPlum,Vej18}
(the second reproducing~\cite{BMPW}); we have not read Goerisch's original papers. A sharper lower
bound is available by refining the mesh ($n=96$ was not run). Other routes would be Liu-type
projection bounds, nonconforming elements in the spirit of~\cite{CG14} (whose interpolation estimate
does not directly apply to projectively glued edges), or an a posteriori bound for the method of
particular solutions.
\end{remark}

\begin{remark}[Strohmaier's mechanism]\label{rem:Str16}
Off $K$, case~(i) of Lemma~\ref{lem:criterion} shows that $\lambda_1$ drops linearly, and since $A(w)$ is
traceless and nonzero the three branches split to first order; this is consistent with the
statement of~\cite{Str16} that no branch is extremal. On $K$ the first-order term vanishes, all three
branches are stationary, and $\lambda_1$ has a strict quadratic maximum. Numerically, along $\FB$ we
observe $\lambda_1=\lambda_1(B)-21.2p^2+O(p^4)$ with $\lambda_2=\lambda_3$ having a minimum, while along the
equilateral family $\FA$ (tangent to $\delta$) $\lambda_1=\lambda_1(B)-5.715|p|+O(p^2)$. These values are
numerical (uncertified finite-element computations) and are not used in the proof. The coincidence
$\lambda_2\equiv\lambda_3$ along $\FB$ is also only observed numerically; the stabiliser of $\wB$ in $G$ has
order $16$ and contains the involution $u\colon x\mapsto i/x$, which suggests an explanation, but we
have not completed it.
\end{remark}

\begin{remark}[Global maximality]
The Main Theorem does not bear on global maximality, which remains open. Combining the rigorous
bound of~\cite{KMP} with Proposition~\ref{prop:LG} shows only that no genus-two surface exceeds
$\lambda_1(B)$ by more than $3.8388976481-3.8388757010$ (about $1.1\times10^{-5}$ with
Hypothesis~\ref{hyp:SU}), and~\cite{Rad} indicates that the existing
single-correlator bootstrap will not close this gap. A quantitative version of our result, giving an
explicit neighbourhood $N$, together with an improved global bound, could in principle reduce the
global problem to a compact region away from $B$. We have not attempted this.
\end{remark}

\begin{remark}[Further questions]
\begin{enumerate}
\item Is $B$ a non-degenerate local maximum in the sharper sense $\lambda_{\min}Q_2|_K<0$, and what is
the second-order behaviour of the individual branches on $K$?
\item The method uses only: a symmetric surface whose $\lambda_1$-cluster is isolated, symmetry forcing
$\tr A\equiv0$, an exact identification of $\ker A$, irreducibility of $\ker A$, and one certified
second-order direction. Fortier Bourque and Petri suggest that $m_1$ should be maximised among local
maximisers of $\lambda_1$~\cite[\S9]{FBP-LP}; it would be natural to apply the same method to other
conjectured maximisers, such as the Klein quartic in genus three~\cite{FBP-K,FBP-LP}.
\item Can the remaining analytic inputs (Theorem~\ref{thm:F1}, Lemma~\ref{lem:N}) and the Arb
certificates be brought into the formal verification?
\end{enumerate}
\end{remark}

\subsection*{Acknowledgements}
The computations, proofs and reviews in this work were produced with Claude (Anthropic) AI agents.
Each component was independently re-checked by separate agent reviewers, who re-ran the key
certificates and re-derived the main formulas. A check by a human expert is pending. The author is
responsible for any remaining errors.

%======================================================================
\appendix
\section{Certificate data}\label{app:data}
%======================================================================

All values are copied from the output files listed in Appendix~\ref{app:repro}. Enclosures
$[m\pm r]$ are Arb balls.

\begin{table}[!htbp]
\caption{Trace formula counts (Lemma~\ref{lem:isolation}(a)); $N=21$.}\label{tab:LP}
\small
\resizebox{\textwidth}{!}{\begin{tabular}{@{}lllll@{}}
\toprule
Run & $s$ & Window & $c\ge$ & Certified bound\\
\midrule
A (Lemma~\ref{lem:LemmaA}) & $1$ & $[0,3.3]$ & $0.9816424759937855$ & $[0.912612352675290375\pm9.95\times10^{-20}]<1$\\
B (\cite[Lemma~3.2]{FBP-K}) & $13/10$ & $[3.3,4.0]$ & $0.9964283521820437$ & $[3.86966111976973759\pm1.65\times10^{-18}]<4$\\
sanity & $13/10$ & $[3.3,3.9]$ & & $3.4803\ \ (\ge3)$\\
sanity & $13/10$ & $[3.3,5.4]$ & & $7.5466\ \ (\ge7)$\\
\bottomrule
\end{tabular}}
\end{table}

\begin{table}[!htbp]
\caption{Certified $P_1$ Rayleigh--Ritz bounds at $B$.}\label{tab:P1}
\small
\begin{tabular}{@{}lll@{}}
\toprule
$n$ & certified $\lambda_3(B)\le$ & area$/4\pi$\\
\midrule
48 & $3.843638188137345$ & $[1.000000\pm1.50\times10^{-8}]$\\
96 & $3.8400767924773938$ & $[1.000000000\pm5.73\times10^{-11}]$\\
192 & $3.839184804702046$ & $[1.00000000000\pm2.23\times10^{-13}]$\\
\bottomrule
\end{tabular}
\end{table}

\begin{table}[!htbp]
\caption{Certified $P_2$ Rayleigh--Ritz bounds at $B$.}\label{tab:P2}
\small
\begin{tabular}{@{}lllll@{}}
\toprule
$n$ & dofs & certified $\lambda_3(B)\le$ & certified $\lambda_1+\lambda_2+\lambda_3\le$ & area$/4\pi$\\
\midrule
24 & 9\,214 & $3.838893606899935$ & $11.516677957238421$ & $[1\pm2.21\times10^{-10}]$\\
48 & 36\,862 & $3.8388876433743824$ & $11.516662767334967$ & $[1\pm1.16\times10^{-14}]$\\
96 & 147\,454 & $3.8388872829261946$ & $11.516661838577997$ & $[1\pm1.19\times10^{-18}]$\\
\bottomrule
\end{tabular}
\end{table}

\begin{table}[!htbp]
\caption{Certificate for $A\ne0$ (per unit $t$, $m(u)$ normalised to $1$). The $n=48$ row fails to
certify.}\label{tab:Anonzero}
\small
\resizebox{\textwidth}{!}{\begin{tabular}{@{}lllllll@{}}
\toprule
$n$ & $\mu=R(u)$ & $b(u,u)$ & $m(w)\le$ & $q(w)\le$ & transfer error & $b(e,e)\le$\\
\midrule
48 & $3.8420704$ & $-4.808678$ & $0.0198$ & $0.0790$ & $9.08$ & $+4.27$\\
96 & $[3.8396833301\pm4.94\times10^{-11}]$ & $[-4.806689\pm1.39\times10^{-7}]$ & $0.00494$ & $0.0198$ & $3.78$ & $-1.0282409843057394$\\
192 & $[3.839086294670\pm5.19\times10^{-13}]$ & $[-4.806192\pm1.96\times10^{-7}]$ & $0.0012$ & $0.0049$ & $1.71$ & $-3.0875694043934345$\\
\midrule
\multicolumn{7}{@{}l}{\emph{SU-free ball (Section~\ref{subsec:transfer}); the rows above use the ball of Hypothesis~\ref{hyp:SU}}}\\
96 & $[3.8396833301\pm4.94\times10^{-11}]$ & $[-4.806689\pm8.00\times10^{-7}]$ & $[0.0050\pm5.95\times10^{-5}]$ & $[0.020\pm2.38\times10^{-4}]$ & $[3.8\pm0.0246]$ & $-0.9989201873540878$\\
192 & $[3.839086294670\pm5.19\times10^{-13}]$ & $[-4.806193\pm8.81\times10^{-7}]$ & $[0.0013\pm6.49\times10^{-5}]$ & $[0.005\pm2.29\times10^{-4}]$ & $[1.7\pm0.0679]$ & $-3.0383476242423058$\\
\bottomrule
\end{tabular}}
\end{table}

\begin{table}[!htbp]
\caption{Sign certificate along $\FB$ ($P_2$).}\label{tab:sign}
\small
\begin{tabular}{@{}lllll@{}}
\toprule
$n$ & dofs & central $g''(0)$ & certified $g''(0)\le$ & transfer error\\
\midrule
24 & 9\,214 & $-14.2738$ & $-9.7138$ & $4.56$\\
48 & 36\,862 & $-14.2742$ & $-12.566084055230016$ & $1.71$\\
96 & 147\,454 & $[-14.2742\pm2.23\times10^{-5}]$ & $-12.850476214196528$ & $1.42$\\
\midrule
96 (SU-free ball) & 147\,454 & $[-14.274\pm2.13\times10^{-4}]$ & $-9.341154560446736$ & \\
\bottomrule
\end{tabular}

\medskip
\parbox{0.9\textwidth}{\footnotesize The first three rows use the ball of Hypothesis~\ref{hyp:SU};
the last row uses the ball of Section~\ref{subsec:transfer} and is the one used in the proof. The rows use the $\pi/2$-symmetry shortcut. Without it:
$n=24$ certifies with the same result, $n=48$ gives $-12.566084126010532$, and an independent re-run
at $n=96$ gives $-12.850476213730866$.}
\end{table}

\begin{table}[!htbp]
\caption{Certified distances to $E\oplus\R1$ at $n=96$ (Lemma~\ref{lem:transfer}).}\label{tab:eta}
\small
\begin{tabular}{@{}llll@{}}
\toprule
$i$ & $\mu_i=R(u_i)$ & $\sqrt{q(w_i)}\le$ & $\sqrt{m(w_i)}\le$\\
\midrule
1 & $[3.838887277826\pm2.34\times10^{-13}]$ & $5.0299\times10^{-3}$ & $2.515\times10^{-3}$\\
2 & $[3.8388872778258\pm8.95\times10^{-14}]$ & $5.0299\times10^{-3}$ & $2.515\times10^{-3}$\\
3 & $[3.838887282926\pm2.30\times10^{-13}]$ & $5.0425\times10^{-3}$ & $2.521\times10^{-3}$\\
\midrule
\multicolumn{4}{@{}l}{\emph{SU-free ball (used in the proof)}}\\
1 & $[3.838887277826\pm2.34\times10^{-13}]$ & $0.016953498968304958$ & $0.008476749477446675$\\
2 & $[3.8388872778258\pm8.95\times10^{-14}]$ & $0.016953498968304958$ & $0.008476749477446675$\\
3 & $[3.838887282926\pm2.30\times10^{-13}]$ & $0.016957233113081952$ & $0.00847861654983665$\\
\bottomrule
\end{tabular}

\medskip
\parbox{0.9\textwidth}{\footnotesize The first three rows use the ball of Hypothesis~\ref{hyp:SU}. There
$\mu_i-\lambda_{\mathrm{low}}\approx1.02\times10^{-6}$ is almost entirely the width of the $\lambda_1$
interval (the discretisation part is about $2\times10^{-8}$). With the wider SU-free ball the distances
are correspondingly larger. In both cases they cannot be reduced by refining the mesh, only by a sharper
lower bound for $\lambda_1$.}
\end{table}

\begin{table}[!htbp]
\caption{Error budget of the sign certificate at $n=96$, for the run with the ball of
Hypothesis~\ref{hyp:SU} (alternative constant; the proof uses $-9.341154560446736$). Rows below the central value are
one-at-a-time contributions (upper end minus central value with only that group switched on); they
are not additive.}\label{tab:budget}
\small
\begin{tabular}{@{}ll@{}}
\toprule
Term & Value\\
\midrule
$T_2=\tr(q''-\ls m'')(u,u)$ (float) & $45.635701833848785$ \\
\quad diagonal (certified) & $[15.613361,\ 15.613361,\ 14.408980]$\\
coupling ($\psi$ terms, float) & $-59.9099$\\
\quad $(q-\ls m)(\psi_i)$ & $0.797,\ 0.797,\ 28.361$\\
central $g''(0)$ & $[-14.2742\pm2.23\times10^{-5}]$\\
\midrule
$X_1$: $4\,b(\psi,w)$ cross terms & $+0.839368$\\
$X_2$: $2(q''-\ls m'')(u,w)$ and $w$--$w$ terms & $+0.570060$\\
$M_0$ & $+4.6\times10^{-4}$\\
$A$, $M_1$, $c_1$ & $1.5\times10^{-5}$, $1.5\times10^{-5}$, $7.4\times10^{-5}$\\
\midrule
certified upper bound (joint ball) & $-12.850476214196528$\\
\bottomrule
\end{tabular}

\medskip
\parbox{0.9\textwidth}{\footnotesize Enclosure maxima at $n=96$: $\beta^1_{\max}=13.2$,
$\beta^2_{\max}=38.1$, $\beta^1_\rho,\beta^2_\rho\le1.49,1.35$; energy-weighted rms values $4.2/9.1$ for
$u$ and $4.6$ ($\psi_1,\psi_2$), $4.0$ ($\psi_3$) for $\psi$.}
\end{table}

\begin{table}[!htbp]
\caption{Rank certificates.}\label{tab:rank}
\small
\begin{tabular}{@{}p{0.42\textwidth}p{0.5\textwidth}@{}}
\toprule
Quantity & Certified value\\
\midrule
$5\times5$ minor of $J$, columns $(0,2,5,7,8)$ & $[46151927.03406868578682884302258\pm2.89\times10^{-24}]$\\
$d(\text{angle sum})/dR$, regular octagon & $[-5.572971195984104378981991909236\pm5.18\times10^{-31}]$\\
$\tr g_k$, $k=0,\dots,3$ & $[22.3137084989848\pm3.97\times10^{-14}]$\\
nonzero $6\times6$ minors of the trace Jacobian & $392$ of $462$\\
best minor, tangent columns $(1,3,5,7,8,10)$ & $[7872428136359854080.0\pm3\times10^{-6}]$\\
$11\times11$ cross-check minor of $[J;J_f]$ & $[-1.3827971011809151404\times10^{24}\pm2.74\times10^{4}]$\\
\bottomrule
\end{tabular}
\end{table}

\begin{table}[!htbp]
\caption{Exact characters (Proposition~\ref{prop:chars} and Section~\ref{sec:symmetry}).}\label{tab:chars}
\small
\begin{tabular}{@{}lcccccc@{}}
\toprule
 & $r$ & $r^2$ & $J$ & $t$ & $\smid$ & $\svert$\\
\midrule
$T_B$ & $-2$ & $-2$ & $6$ & $0$ & $0$ & $0$\\
$K_0$ & $-1$ & & & $0$ & $1$ & $-1$\\
$W_0$ & $-1$ & & & $0$ & $-1$ & $1$\\
\bottomrule
\end{tabular}

\medskip
\parbox{0.9\textwidth}{\footnotesize $\langle\chi_T,\chi_T\rangle=2$, $\langle\chi_T,1\rangle=0$;
$K_0$, $W_0$ have norm $1$ and are orthogonal; $\chi_T$ vanishes on all orientation-reversing
elements. Group: order $96$; $G^+$ element orders $1{:}1$, $2{:}13$, $3{:}8$, $4{:}6$, $6{:}8$, $8{:}12$;
$D_8$ contains $4$ reflections of $\smid$-type and $4$ of $\svert$-type.}
\end{table}

\clearpage
%======================================================================
\section{Reproducibility}\label{app:repro}
%======================================================================

The scripts are in \texttt{work/} of the project folder
(\texttt{items/23-bolza-lambda1-extremality/}) and are run from that directory with the virtual
environment \texttt{work/.venv} (python-flint~0.9.0, sympy~1.14.0, numpy, scipy), under
\texttt{nice -n 10} and \texttt{OMP\_NUM\_THREADS=2}. Times are wall-clock on a shared laptop.

\begin{center}
\small
\begin{tabular}{@{}p{0.47\textwidth}p{0.3\textwidth}p{0.13\textwidth}@{}}
\toprule
Command & Output & Time\\
\midrule
\texttt{python lp\_count.py A 3.3 21 1.0} & \texttt{lp\_count.out}, \texttt{lp\_A\_3.3.json} & $\sim3$ s\\
\texttt{python lp\_count.py B 3.3 4.0 21 1.3} & \texttt{lp\_B\_3.3\_4.0.json} & $\sim3$ s\\
\texttt{python cert\_rr.py 48 96 192} & \texttt{cert\_rr\_bolza.out} & $\sim8$ min\\
\texttt{python cert\_rr\_p2.py 24 48 96} & \texttt{cert\_rr\_p2\_bolza.out} & $\sim100$ s\\
\texttt{SYM=1 python cert\_first\_order.py 96 192} & \texttt{cert\_first\_order.out} & $\sim12$ min\\
\texttt{python sym\_check\_elements.py 16} & \texttt{sym\_check\_elements.out} & $\sim1$ min\\
\texttt{python cert\_rank.py} & \texttt{cert\_rank.out} & $\sim5$ s\\
\texttt{python cert\_dpi\_rank.py} & \texttt{cert\_dpi\_rank.out}, \texttt{.json} & $<1$ s\\
\texttt{python cert\_chars.py} & \texttt{cert\_chars\_A.out}, \texttt{.json} & $\sim22$ s\\
\texttt{python cert\_chars.py B 48 96} (not used in the proof) & \texttt{cert\_chars\_B.out} & $\sim4.4$ min\\
\texttt{python cert\_sign.py 48 96} & \texttt{cert\_sign.out}, \texttt{cert\_sign\_48\_96.json} & $\sim4.3$ min\\
\texttt{SYM=0 python cert\_sign.py 12 24} & \texttt{cert\_sign\_nosym.out} & \\
\texttt{SYM=0 python cert\_sign.py 48} & \texttt{cert\_sign\_48\_nosym.out} & \\
\texttt{python lb\_torus.py} (validation) & \texttt{lb\_torus.out} & \\
\texttt{python lb\_bolza.py 48} & \texttt{lb\_bolza\_48.out}, \texttt{.json} & $211$ s\\
\texttt{python lb\_cert\_sign.py 96} & \texttt{lb\_cert\_sign.out}, \texttt{lb\_cert\_sign\_96.json} & $165$ s\\
\texttt{python lb\_cert\_first\_order.py 96 192} & \texttt{lb\_cert\_first\_order.out}, \texttt{lb\_cert\_first\_order\_}\allowbreak\texttt{96\_192.json} & $195$ s $+$ $718$ s\\
\bottomrule
\end{tabular}
\end{center}

The core floating-point geometry is in \texttt{octagon.py}. The Lean development is in
\texttt{lean/}; it is built by \texttt{lean/build.sh}, which compiles the files
\texttt{RepTheory}, \texttt{LocalMax}, \texttt{Taylor}, \texttt{RepReal}, \texttt{Certificates},
\texttt{KyFan}, \texttt{Assembly} and \texttt{Model} with plain \texttt{lean} against prebuilt Mathlib
\texttt{.olean} files and then prints the axioms of the main declarations (a clean build took
2~min~33~s).

\emph{Process.} The certificates and written proofs were produced in an AI-assisted workflow
(see the acknowledgements). Each component was re-checked by a separate reviewer: the
trace-formula, $P_1$, $\tr A$ and $A\ne0$ certificates; the sign certificate (including an
independent re-derivation of Lemma~\ref{lem:g2formula}, an independent $P_1$ code, sensitivity
analysis and negative controls); the representation theory (with an independently written model);
the proof of Lemma~\ref{lem:N}; the chart certificate and Lemma~\ref{lem:Ginv}, whose
certificate was re-run with bit-identical output; and the Lehmann--Goerisch lower bound
(Proposition~\ref{prop:LG}) with the re-certified first- and second-order certificates. That
reviewer checked the quoted theorem against both secondary sources, redid the final step in
independent interval arithmetic, re-ran the $n=24$ and torus cases, and re-ran both downstream
certificates. The review records are kept with the project files.

\bibliographystyle{amsplain}
\bibliography{refs}

\end{document}
